\input{preamble}

% OK, start here.
%
\begin{document}

\title{Compléments sur les groupoïdes en espaces algébriques}


\maketitle

\phantomsection
\label{section-phantom}

\tableofcontents

\section{Introduction}
\label{section-introduction}

\noindent
Ce chapitre est consacré à des questions avancées sur les groupoïdes
en espaces algébriques.
Bien que les résultats soient énoncés en termes de groupoïdes en
espaces algébriques, le
lecteur doit garder à l'esprit le diagramme $2$-cartésien
\begin{equation}
\label{equation-quotient-stack}
\vcenter{
\xymatrix{
R \ar[r] \ar[d] & U \ar[d] \\
U \ar[r] & [U/R]
}
}
\end{equation}
où $[U/R]$ est le champ quotient; voir
Groupoïdes dans les espaces, remarque \ref{spaces-groupoids-remark-fundamental-square}.
Nombre des résultats sont motivés par l'étude de ce diagramme.
Voir par exemple le bel article \cite{K-M} de Keel et Mori.





\section{Notations}
\label{section-notation}

\noindent
Nous continuons à suivre les conventions et notations introduites dans
Groupoïdes dans les espaces, section \ref{spaces-groupoids-section-notation}.





\section{Diagrammes utiles}
\label{section-diagrams}

\noindent
Rappelons brièvement les résultats de
Groupoïdes dans les espaces, lemmes \ref{spaces-groupoids-lemma-diagram} et
\ref{spaces-groupoids-lemma-diagram-pull}
afin de pouvoir nous y référer aisément dans ce chapitre.
Soit $S$ un schéma. Soit $B$ un espace algébrique au-dessus de $S$.
Soit $(U, R, s, t, c)$ un groupoïde en espaces algébriques au-dessus de $B$.
Dans le diagramme commutatif
\begin{equation}
\label{equation-diagram}
\vcenter{
\xymatrix{
& U & \\
R \ar[d]_s \ar[ru]^t &
R \times_{s, U, t} R
\ar[l]^-{\text{pr}_0} \ar[d]^{\text{pr}_1} \ar[r]_-c &
R \ar[d]^s \ar[lu]_t \\
U & R \ar[l]_t \ar[r]^s & U
}
}
\end{equation}
les deux carrés inférieurs sont cartésiens.
De plus, le triangle supérieur (qui est en réalité un carré)
est lui aussi cartésien.

\medskip\noindent
Le diagramme
\begin{equation}
\label{equation-pull}
\vcenter{
\xymatrix{
R \times_{t, U, t} R
\ar@<1ex>[r]^-{\text{pr}_1} \ar@<-1ex>[r]_-{\text{pr}_0}
\ar[d]_{\text{pr}_0 \times c \circ (i, 1)} &
R \ar[r]^t \ar[d]^{\text{id}_R} &
U \ar[d]^{\text{id}_U} \\
R \times_{s, U, t} R
\ar@<1ex>[r]^-c \ar@<-1ex>[r]_-{\text{pr}_0} \ar[d]_{\text{pr}_1} &
R \ar[r]^t \ar[d]^s &
U \\
R \ar@<1ex>[r]^s \ar@<-1ex>[r]_t &
U
}
}
\end{equation}
est commutatif. Les deux lignes supérieures sont isomorphes par les flèches
verticales indiquées. Les deux carrés inférieurs de gauche sont cartésiens.






\section{Structure locale}
\label{section-local}

\noindent
Soit $S$ un schéma.
Soit $(U, R, s, t, c, e, i)$ un groupoïde en espaces algébriques au-dessus de $S$.
Soit $\overline{u}$ un point géométrique de $U$. Dans cette section, nous
décrivons la structure obtenue sur les anneaux locaux
(Propriétés des espaces, définition
\ref{spaces-properties-definition-etale-local-rings})
$$
A = \mathcal{O}_{U, \overline{u}}
\quad\text{et}\quad
B = \mathcal{O}_{R, e(\overline{u})}
$$
Nous convenons de désigner par les lettres correspondantes les homomorphismes
d'anneaux locaux induits par les morphismes $s, t, c, e, i$.
En particulier, nous avons un diagramme commutatif
$$
\xymatrix{
A \ar[rd]_t \ar[rrd]^1 \\
& B \ar[r]^e & A \\
A \ar[ru]^s \ar[rru]_1
}
$$
d'anneaux locaux. Ainsi, si $I \subset B$ désigne le noyau de $e : B \to A$,
alors $B = s(A) \oplus I = t(A) \oplus I$. Posons
$$
C = \mathcal{O}_{R \times_{s, U, t} R, (e, e)(\overline{u})}
$$
Alors
$$
C =
(B \otimes_{s, A, t} B)_{\mathfrak m_B \otimes B + B \otimes \mathfrak m_B}^h
$$
car le localisé
$(B \otimes_{s, A, t} B)_{\mathfrak m_B \otimes B + B \otimes \mathfrak m_B}$
a un corps résiduel séparablement clos.
Soit $J \subset C$ l'idéal de $C$ engendré par $I \otimes B + B \otimes I$.
Alors $J$ est aussi le noyau de l'homomorphisme d'anneaux locaux
$$
(e, e) : C \longrightarrow A
$$
La loi de composition $c : R \times_{s, U, t} R \to R$ correspond à un
homomorphisme d'anneaux
$$
c : B \longrightarrow C
$$
qui envoie $I$ dans $J$.

\begin{lemma}
\label{lemma-first-order-structure-c}
L'application $I/I^2 \to J/J^2$ induite par $c$ est la composée
$$
I/I^2 \xrightarrow{(1, 1)} I/I^2 \oplus I/I^2 \to J/J^2
$$
où la seconde flèche provient de l'égalité
$J = (I \otimes B + B \otimes I)C$.
L'application $i : B \to B$ induit l'application $-1 : I/I^2 \to I/I^2$.
\end{lemma}

\begin{proof}
Pour décrire un homomorphisme local de $C$ vers un autre
anneau local hensélien, il suffit de préciser l'image
des éléments de la forme $b_1 \otimes b_2$, par exemple d'après
Algèbre, lemme \ref{algebra-lemma-henselian-functorial}.
Compte tenu de cela, nous avons les deux applications canoniques
$$
e_2 : C \to B,\ b_1 \otimes b_2 \mapsto b_1s(e(b_2)),\quad
e_1 : C \to B,\ b_1 \otimes b_2 \mapsto t(e(b_1))b_2
$$
correspondant aux immersions
$R \to R \times_{s, U, t} R$ données par
$r \mapsto (r, e(s(r)))$ et $r \mapsto (e(t(r)), r)$.
Ces applications définissent des applications $J/J^2 \to I/I^2$ qui,
prises ensemble, fournissent un inverse de l'application
$I/I^2 \oplus I/I^2 \to J/J^2$ du lemme. Pour démontrer l'assertion, il suffit
donc de montrer que $e_1 \circ c : B \to B$ et $e_2 \circ c : B \to B$
sont les applications identiques. Cela résulte du fait que les deux
composées $R \to R \times_{s, U, t} R \to R$ sont des identités.

\medskip\noindent
L'assertion concernant $i$ résulte de celle concernant $c$ et du
fait que $c \circ (1, i) = e \circ t$. Certains détails sont omis.
\end{proof}








\section{Groupoïde des sections}
\label{section-groupoid-sections}

\noindent
Supposons donné un groupoïde $(\text{Ob}, \text{Flèches}, s, t, c, e, i)$.
On peut alors construire un monoïde $\Gamma$ dont les éléments sont les
applications $\delta : \text{Ob} \to \text{Flèches}$ telles que
$s \circ \delta = \text{id}_{\text{Ob}}$, la composition étant donnée par
$$
\delta_1 \circ \delta_2 = c(\delta_1 \circ t \circ \delta_2, \delta_2)
$$
Autrement dit, un élément de $\Gamma$ est une règle $\delta$ qui
prescrit une flèche issue de chaque objet, et la composition
est celle que l'on attend naturellement. Par exemple,
$$
\vcenter{
\xymatrix{
& \bullet \ar[dl] \\
\bullet \ar@(ul, dl)[] & \bullet \ar[d] \\
& \bullet \ar[lu]
}
}
\quad\circ\quad
\vcenter{
\xymatrix{
& \bullet \ar[d] \\
\bullet \ar[ru] & \bullet \ar[d] \\
& \bullet \ar[lu]
}
}
\quad = \quad\quad
\vcenter{
\xymatrix{
& \bullet \ar@/^/[dd] \\
\bullet \ar@(ul, dl)[] & \bullet \ar[l] \\
& \bullet \ar[lu]
}
}
$$
avec des notations évidentes.

\medskip\noindent
Le même procédé s'applique à un groupoïde en espaces
algébriques $(U, R, s, t, c, e, i)$ au-dessus d'un schéma $S$.
Plus précisément, on prend pour éléments de $\Gamma$ l'ensemble
$$
\Gamma = \{\delta : U \to R \mid s \circ \delta = \text{id}_U\}
$$
et la composition $\circ : \Gamma \times \Gamma \to \Gamma$
est donnée par la règle ci-dessus :
\begin{equation}
\label{equation-composition}
\delta_1 \circ \delta_2 = c(\delta_1 \circ t \circ \delta_2, \delta_2)
\end{equation}
L'élément neutre est $e \in \Gamma$. En général, le monoïde $\Gamma$ n'est pas
un groupe, car certains éléments $\delta \in \Gamma$ peuvent ne pas
avoir d'inverse. En effet, il est clair que $\delta \in \Gamma$
a un inverse si et seulement si $t \circ \delta$ est un automorphisme
de $U$; dans ce cas,
$\delta^{-1} = i \circ \delta \circ (t \circ \delta)^{-1}$.

\medskip\noindent
Pour un usage ultérieur, examinons le sous-groupe
$\Gamma_0$ de $\Gamma$ formé des sections infinitésimalement
proches de l'identité $e$. Plus précisément, supposons donné un
sous-espace fermé $R$-invariant $U_0 \subset U$ tel que $U$
soit un épaississement du premier ordre de $U_0$. Notons
$R_0 = s^{-1}(U_0) = t^{-1}(U_0)$ et soit
$(U_0, R_0, s_0, t_0, c_0, e_0, i_0)$ le groupoïde
en espaces algébriques correspondant. Posons
$$
\Gamma_0 = \{\delta \in \Gamma \mid \delta|_{U_0} = e_0\}
$$
Si $s$ et $t$ sont plats, tout élément de $\Gamma_0$
est inversible. En effet, $t \circ \delta$ est alors un
morphisme $U \to U$ qui induit l'identité sur
$\mathcal{O}_{U_0}$ et sur $\mathcal{C}_{U_0/U}$
(lemme \ref{lemma-idenity-on-conormal});
on conclut grâce à la suite exacte
courte
$0 \to \mathcal{C}_{U_0/U} \to \mathcal{O}_U \to \mathcal{O}_{U_0} \to 0$.

\begin{lemma}
\label{lemma-idenity-on-conormal}
Dans la situation étudiée dans cette section, soient $\delta \in \Gamma_0$
et $f = t \circ \delta : U \to U$. Si $s, t$ sont plats, alors l'application
canonique $\mathcal{C}_{U_0/U} \to \mathcal{C}_{U_0/U}$ induite par $f$
(Compléments sur les morphismes d'espaces, lemme
\ref{spaces-more-morphisms-lemma-conormal-functorial})
est l'application identique.
\end{lemma}

\begin{proof}
Pour le voir, prolongeons comme suit la partie inférieure du diagramme
(\ref{equation-pull}) :
$$
\xymatrix{
Y \ar[r] \ar[d] &
R \times_{s, U, t} R
\ar@<1ex>[r]^-c \ar@<-1ex>[r]_-{\text{pr}_0} \ar[d]_{\text{pr}_1} &
R \ar[r]^t \ar[d]^s &
U \\
U \ar[r]_\delta &
R \ar@<1ex>[r]^s \ar@<-1ex>[r]_t &
U
}
$$
où le carré de gauche est cartésien; ceci constitue notre définition
de $Y$; nous n'aurons pas besoin d'en savoir davantage sur $Y$.
Un changement de base partout à $U_0$ donne un diagramme analogue
ayant les mêmes propriétés.
Nous voulons montrer que $\text{id}_U = s \circ \delta$
et $f = t \circ \delta$ induisent les mêmes applications sur les faisceaux
conormaux. Comme $s$ est plat et surjectif, il suffit de démontrer la même
chose pour les deux composées $a, b : Y \to R$ de la ligne supérieure.
Observons que $a_0 = b_0$ et que l'un de $a$ et $b$ est un isomorphisme,
puisque nous savons que $s \circ \delta$ est un isomorphisme. Par conséquent,
les deux morphismes $a, b : Y \to R$ sont des morphismes entre espaces
algébriques plats au-dessus de $U$ (par le morphisme $t : R \to U$
et le morphisme $t \circ a = t \circ b : Y \to U$).
Cela donne le résultat voulu. En effet, par compatibilité aux compositions dans
Compléments sur les morphismes d'espaces, lemme
\ref{spaces-more-morphisms-lemma-conormal-functorial-more}
on conclut que les deux applications
$a_0^*\mathcal{C}_{R_0/R} \to \mathcal{C}_{Y_0/Y}$
s'insèrent dans un diagramme commutatif
$$
\xymatrix{
a_0^*\mathcal{C}_{R_0/R} \ar[rr] & & \mathcal{C}_{Y_0/Y} \\
a_0^*t_0^*\mathcal{C}_{U_0/U} \ar[u] \ar@{=}[rr] & &
(t_0 \circ a_0)^*\mathcal{C}_{U_0/U} \ar[u]
}
$$
dont les flèches verticales sont des isomorphismes d'après
Compléments sur les morphismes d'espaces, lemme
\ref{spaces-more-morphisms-lemma-deform}.
Le lemme en résulte.
\end{proof}

\noindent
Identifions maintenant le groupe $\Gamma_0$.
Appliquons la discussion de Compléments sur les morphismes d'espaces, remarques
\ref{spaces-more-morphisms-remark-action-by-derivations} et
\ref{spaces-more-morphisms-remark-another-special-case}
au diagramme
$$
\xymatrix{
(U_0 \subset U) \ar@{..>}[rr]_{(e_0, \delta)}
\ar[rd]_{(\text{id}_{U_0}, \text{id}_U)} & &
(R_0 \subset R) \ar[ld]^{(s_0, s)} \\
& (U_0 \subset U)
}
$$
On voit que $\delta = \theta \cdot e$ pour une unique application
$\mathcal{O}_{U_0}$-linéaire
$\theta : e_0^*\Omega_{R_0/U_0} \to \mathcal{C}_{U_0/U}$.
On obtient ainsi une bijection
\begin{equation}
\label{equation-isomorphism}
\Hom_{\mathcal{O}_{U_0}}(e_0^*\Omega_{R_0/U_0}, \mathcal{C}_{U_0/U})
\longrightarrow
\Gamma_0
\end{equation}
en appliquant Compléments sur les morphismes d'espaces, lemme
\ref{spaces-more-morphisms-lemma-action-sheaf}.

\begin{lemma}
\label{lemma-composition-is-addition}
La bijection (\ref{equation-isomorphism}) est un isomorphisme
de groupes.
\end{lemma}

\begin{proof}
Soient $\delta_1, \delta_2 \in \Gamma_0$ correspondant à $\theta_1, \theta_2$
comme ci-dessus, et la composée $\delta = \delta_1 \circ \delta_2$
dans $\Gamma_0$ correspondant à $\theta$. Il faut montrer que
$\theta = \theta_1 + \theta_2$. Rappelons
(Compléments sur les morphismes d'espaces, lemme
\ref{spaces-more-morphisms-lemma-action-by-derivations})
que $\theta_1, \theta_2, \theta$ correspondent aux dérivations
$D_1, D_2, D : e_0^{-1}\mathcal{O}_{R_0} \to \mathcal{C}_{U_0/U}$
données par $D_1 = \theta_1 \circ \text{d}_{R_0/U_0}$, et ainsi de suite.
Il suffit de vérifier que $D = D_1 + D_2$.

\medskip\noindent
On peut vérifier l'égalité sur les fibres.
Soit $\overline{u}$ un point géométrique de $U$ et utilisons les anneaux
locaux $A, B, C$ introduits dans la section \ref{section-local}.
Aux morphismes $\delta_i$ correspondent des homomorphismes d'anneaux
$\delta_i : B \to A$. Soit $K \subset A$ l'idéal de carré nul tel que
$A/K = \mathcal{O}_{U_0, \overline{u}}$.
Autrement dit, $K$ est la fibre de $\mathcal{C}_{U_0/U}$ en $\overline{u}$.
La condition $\delta_i \in \Gamma_0$
signifie exactement que $\delta_i(I) \subset K$.
La dérivation $D_i$ est simplement l'application $\delta_i - e : B \to A$.
Comme $B = s(A) \oplus I$, la dérivation $D_i$ est déterminée par
sa restriction à $I$, laquelle est précisément donnée par
$\delta_i|_I$. De plus, $D_i$, et donc $\delta_i$, annule $I^2$
puisque $\delta_i(I) \subset K$ et $K^2 = 0$.

\medskip\noindent
Pour achever la démonstration, observons que $\delta$ correspond à la
composée
$$
B \to C =
(B \otimes_{s, A, t} B)^h_{\mathfrak m_B \otimes B + B \otimes \mathfrak m_B}
\to A
$$
où la première flèche est $c$ et la seconde est déterminée
par la règle
$b_1 \otimes b_2 \mapsto \delta_2(t(\delta_1(b_1))) \delta_2(b_2)$
comme il résulte de (\ref{equation-composition}).
D'après le lemme \ref{lemma-first-order-structure-c},
on voit qu'un élément $\zeta$ de $I$ est envoyé sur
$\zeta \otimes 1 + 1 \otimes \zeta$, à des termes d'ordre supérieur près.
On en conclut que
$$
D(\zeta) = (\delta_2 \circ t)\left(D_1(\zeta)\right) + D_2(\zeta)
$$
Or, d'après le lemme \ref{lemma-idenity-on-conormal},
l'action de $\delta_2 \circ t$ sur
$K = \mathcal{C}_{U_0/U, \overline{u}}$ est l'identité, ce qui
achève la preuve.
\end{proof}






\section{Propriétés des groupoïdes}
\label{section-technical-lemma}

\noindent
Cette section est l'analogue de
Compléments sur les groupoïdes, section \ref{more-groupoids-section-technical-lemma}.
Il est vivement conseillé au lecteur de commencer par lire cette section.

\medskip\noindent
Le lemme suivant est l'analogue de
Compléments sur les groupoïdes, lemme \ref{more-groupoids-lemma-property-invariant}.

\begin{lemma}
\label{lemma-property-invariant}
Soit $B \to S$ comme dans la section \ref{section-notation}.
Soit $(U, R, s, t, c)$ un groupoïde en espaces algébriques au-dessus de $B$.
Soit
$\tau \in \{fppf, \linebreak[0] \etale, \linebreak[0]
lisse, \linebreak[0] syntomique\}$.
Soit $\mathcal{P}$ une propriété des morphismes d'espaces algébriques
qui est $\tau$-locale sur le but
(Descente sur les espaces,
définition \ref{spaces-descent-definition-property-morphisms-local}).
Supposons que $\{s : R \to U\}$ et $\{t : R \to U\}$ soient des recouvrements
pour la topologie $\tau$. Soit $W \subset U$ le plus grand sous-espace ouvert
tel que $s^{-1}(W) \to W$ possède la propriété $\mathcal{P}$.
Alors $W$ est $R$-invariant
(Groupoïdes dans les espaces,
définition \ref{spaces-groupoids-definition-invariant-open}).
\end{lemma}

\begin{proof}
L'existence et les propriétés de l'ouvert $W \subset U$ sont décrites dans
Descente sur les espaces, lemme \ref{spaces-descent-lemma-largest-open-of-the-base}.
Dans le
diagramme (\ref{equation-diagram}),
soit $W_1 \subset R$ le plus grand sous-schéma ouvert sur lequel le morphisme
$\text{pr}_1 : R \times_{s, U, t} R \to R$ possède la propriété $\mathcal{P}$.
Il résulte du lemme précité de
Descente sur les espaces, lemme \ref{spaces-descent-lemma-largest-open-of-the-base},
et de l'hypothèse que $\{s : R \to U\}$ et $\{t : R \to U\}$ sont des
recouvrements pour la topologie $\tau$, que $t^{-1}(W)=W_1=s^{-1}(W)$, comme voulu.
\end{proof}

\begin{lemma}
\label{lemma-property-G-invariant}
Soit $B \to S$ comme dans la section \ref{section-notation}.
Soit $(U, R, s, t, c)$ un groupoïde en espaces algébriques au-dessus de $B$.
Soit $G \to U$ son espace algébrique en groupes stabilisateur.
Soit
$\tau \in \{fppf, \linebreak[0] \etale, \linebreak[0]
lisse, \linebreak[0] syntomique\}$.
Soit $\mathcal{P}$ une propriété des morphismes d'espaces algébriques
qui est $\tau$-locale sur le but.
Supposons que $\{s : R \to U\}$ et $\{t : R \to U\}$ soient des recouvrements
pour la topologie $\tau$. Soit $W \subset U$ le plus grand sous-espace ouvert
tel que $G_W \to W$ possède la propriété $\mathcal{P}$.
Alors $W$ est $R$-invariant (voir
Groupoïdes dans les espaces,
définition \ref{spaces-groupoids-definition-invariant-open}).
\end{lemma}

\begin{proof}
L'existence et les propriétés de l'ouvert $W \subset U$ sont décrites dans
Descente sur les espaces, lemme \ref{spaces-descent-lemma-largest-open-of-the-base}.
Le morphisme
$$
G \times_{U, t} R \longrightarrow R \times_{s, U} G, \quad
(g, r) \longmapsto (r, r^{-1} \circ g \circ r)
$$
est un isomorphisme d'espaces algébriques au-dessus de $R$ (où $\circ$ désigne
la composition dans le groupoïde). Ainsi $s^{-1}(W)=t^{-1}(W)$ d'après les
propriétés de $W$ démontrées dans le lemme précité de
Descente sur les espaces, lemme \ref{spaces-descent-lemma-largest-open-of-the-base}.
\end{proof}




\section{Comparaison des fibres}
\label{section-fibres}

\noindent
Cette section est l'analogue de
Compléments sur les groupoïdes, section \ref{more-groupoids-section-fibres}.
Il est vivement conseillé au lecteur de commencer par lire cette section.

\begin{lemma}
\label{lemma-two-fibres}
Soit $B \to S$ comme dans la section \ref{section-notation}.
Soit $(U, R, s, t, c)$ un groupoïde en espaces algébriques au-dessus de $B$.
Soit $K$ un corps et soient $r, r' : \Spec(K) \to R$
des morphismes tels que $t \circ r = t \circ r' : \Spec(K) \to U$.
Posons $u = s \circ r$, $u' = s \circ r'$ et notons
$F_u = \Spec(K) \times_{u, U, s} R$ et
$F_{u'} = \Spec(K) \times_{u', U, s} R$ les produits fibrés.
Alors $F_u \cong F_{u'}$ comme espaces algébriques au-dessus de $K$.
\end{lemma}

\begin{proof}
Nous utilisons l'existence et les propriétés du
diagramme (\ref{equation-diagram}).
Il existe un morphisme $\xi : \Spec(K) \to R \times_{s, U, t} R$
tel que $\text{pr}_0 \circ \xi = r$ et $c \circ \xi = r'$.
Soit $\tilde r = \text{pr}_1 \circ \xi : \Spec(K) \to R$.
En examinant alors les deux carrés inférieurs du
diagramme (\ref{equation-diagram}),
on voit que $F_u$ et $F_{u'}$ s'identifient tous deux à l'espace algébrique
$\Spec(K) \times_{\tilde r, R, \text{pr}_1} (R \times_{s, U, t} R)$.
\end{proof}

\noindent
En fait, dans la situation du lemme, les morphismes de paires
$s : (R, r) \to (U, u)$ et $s : (R, r') \to (U, u')$ sont
localement isomorphes pour la topologie $\tau$, pourvu que $\{s: R \to U\}$
soit un recouvrement $\tau$. Nous insérerons ici un énoncé précis au besoin.








\section{Restriction des groupoïdes}
\label{section-restricting-groupoids}

\noindent
Dans cette section, nous rassemblons plusieurs lemmes sur les
propriétés des groupoïdes qui se transmettent aux restrictions.
La plupart se démontrent en considérant le diagramme
définissant la restriction :
\begin{equation}
\label{equation-restriction}
\vcenter{
\xymatrix{
R' \ar[d] \ar[r] \ar@/_3pc/[dd]_{t'} \ar@/^1pc/[rr]^{s'}&
R \times_{s, U} U' \ar[r] \ar[d] &
U' \ar[d]^g \\
U' \times_{U, t} R \ar[d] \ar[r] &
R \ar[r]^s \ar[d]_t &
U \\
U' \ar[r]^g &
U
}
}
\end{equation}
Voir
Groupoïdes dans les espaces, lemme \ref{spaces-groupoids-lemma-restrict-groupoid}.

\begin{lemma}
\label{lemma-restrict-preserves-type}
Soit $S$ un schéma. Soit $B$ un espace algébrique au-dessus de $S$.
Soit $(U, R, s, t, c)$ un groupoïde en espaces algébriques au-dessus de $B$.
Soit $g : U' \to U$ un morphisme d'espaces algébriques au-dessus de $B$.
Soit $(U', R', s', t', c')$ la restriction de
$(U, R, s, t, c)$ par $g$.
\begin{enumerate}
\item Si $s, t$ sont localement de type fini et $g$ est localement de type
fini, alors $s', t'$ sont localement de type fini.
\item Si $s, t$ sont localement de présentation finie et $g$ est localement de
présentation finie, alors $s', t'$ sont localement de présentation finie.
\item Si $s, t$ sont plats et $g$ est plat, alors $s', t'$ sont plats.
\item Ajouter d'autres cas ici.
\end{enumerate}
\end{lemma}

\begin{proof}
La propriété d'être localement de type fini est stable par composition
et par changement de base arbitraire; voir
Morphismes d'espaces,
lemmes \ref{spaces-morphisms-lemma-composition-finite-type} et
\ref{spaces-morphisms-lemma-base-change-finite-type}.
Le point (1) résulte donc du diagramme (\ref{equation-restriction}).
Pour les autres cas, voir
Morphismes d'espaces,
lemmes \ref{spaces-morphisms-lemma-composition-finite-presentation},
\ref{spaces-morphisms-lemma-base-change-finite-presentation},
\ref{spaces-morphisms-lemma-composition-flat} et
\ref{spaces-morphisms-lemma-base-change-flat}.
\end{proof}









\section{Propriétés des groupes sur les corps et des groupoïdes sur les corps}
\label{section-properties-groupoids-on-fields}

\noindent
Il est conseillé au lecteur de consulter d'abord les sections correspondantes
pour les schémas en groupoïdes; voir
Groupoïdes, section \ref{groupoids-section-properties-group-schemes-field}
et
Compléments sur les groupoïdes,
section \ref{more-groupoids-section-properties-groupoids-on-fields}.

\begin{situation}
\label{situation-group-over-field}
Ici, $S$ est un schéma, $k$ est un corps au-dessus de $S$, et
$(G, m)$ est un espace algébrique en groupes au-dessus de $\Spec(k)$.
\end{situation}

\begin{situation}
\label{situation-groupoid-on-field}
Ici, $S$ est un schéma, $B$ est un espace algébrique, et
$(U, R, s, t, c)$ est un groupoïde en espaces algébriques au-dessus de $B$
avec $U = \Spec(k)$ pour un certain corps $k$.
\end{situation}

\noindent
Remarquons que, dans la
situation \ref{situation-group-over-field},
on obtient un groupoïde en espaces algébriques
\begin{equation}
\label{equation-groupoid-from-group}
(\Spec(k), G, p, p, m)
\end{equation}
où $p : G \to \Spec(k)$ est le morphisme structural de $G$; voir
Groupoïdes dans les espaces, lemme \ref{spaces-groupoids-lemma-groupoid-from-action}.
On se trouve alors dans la
situation \ref{situation-groupoid-on-field}.
Nous l'utiliserons sans autre mention dans la suite de cette section.

\begin{lemma}
\label{lemma-groupoid-on-field-open-multiplication}
Dans la
situation \ref{situation-groupoid-on-field},
le morphisme de composition $c : R \times_{s, U, t} R \to R$ est plat et
universellement ouvert.
Dans la
situation \ref{situation-group-over-field},
la loi de groupe $m : G \times_k G \to G$ est plate et
universellement ouverte.
\end{lemma}

\begin{proof}
La composition est isomorphe à la projection
$\text{pr}_1 : R \times_{t, U, t} R \to R$ d'après le
diagramme (\ref{equation-pull}).
La projection est plate, comme changement de base du morphisme plat $t$,
et ouverte d'après
Morphismes d'espaces,
lemme \ref{spaces-morphisms-lemma-space-over-field-universally-open}.
La seconde assertion résulte aussitôt de la première, puisque
$m$ correspond à $c$ dans (\ref{equation-groupoid-from-group}).
\end{proof}

\noindent
Remarquons que le lemme suivant s'applique notamment lorsqu'on travaille
avec des espaces algébriques quasi-séparés ou localement séparés
(Espaces décents, lemme \ref{decent-spaces-lemma-locally-separated-decent}).

\begin{lemma}
\label{lemma-group-scheme-over-field-separated}
Dans la situation \ref{situation-groupoid-on-field},
supposons que $R$ soit un espace décent. Alors $R$ est un espace algébrique séparé.
Dans la situation \ref{situation-group-over-field}, supposons que
$G$ soit un espace algébrique décent. Alors $G$ est un espace algébrique séparé.
\end{lemma}

\begin{proof}
Commençons par démontrer la seconde assertion. D'après Groupoïdes dans les espaces,
lemme \ref{spaces-groupoids-lemma-group-scheme-separated},
il faut montrer que $e : S \to G$ est une immersion fermée.
Cela résulte d'Espaces décents, lemme
\ref{decent-spaces-lemma-finite-residue-field-extension-finite}.

\medskip\noindent
Démontrons ensuite la première assertion. On peut pour cela remplacer $B$ par $S$.
D'après le paragraphe précédent, le schéma en groupes stabilisateur $G \to U$ est séparé. D'après
Groupoïdes dans les espaces, lemme \ref{spaces-groupoids-lemma-diagonal},
le morphisme $j = (t, s) : R \to U \times_S U$ est séparé.
Comme $U$ est le spectre d'un corps, le schéma
$U \times_S U$ est affine (par la construction des produits fibrés dans
Schémas, section \ref{schemes-section-fibre-products}).
Par conséquent, $R$ est séparé; voir
Morphismes d'espaces, lemme
\ref{spaces-morphisms-lemma-separated-over-separated}.
\end{proof}

\begin{lemma}
\label{lemma-restrict-groupoid-on-field}
Dans la
situation \ref{situation-groupoid-on-field}.
Soit $k'/k$ une extension de corps, $U' = \Spec(k')$,
et soit $(U', R', s', t', c')$ la restriction de
$(U, R, s, t, c)$ par $U' \to U$. Dans le diagramme de définition
$$
\xymatrix{
R' \ar[d] \ar[r] \ar@/_3pc/[dd]_{t'} \ar@/^1pc/[rr]^{s'} \ar@{..>}[rd] &
R \times_{s, U} U' \ar[r] \ar[d] &
U' \ar[d] \\
U' \times_{U, t} R \ar[d] \ar[r] &
R \ar[r]^s \ar[d]_t &
U \\
U' \ar[r] &
U
}
$$
tous les morphismes sont surjectifs, plats et universellement ouverts.
La flèche pointillée $R' \to R$ est de plus affine.
\end{lemma}

\begin{proof}
Le morphisme $U' \to U$ est $\Spec(k') \to \Spec(k)$;
il est donc affine, surjectif et plat. Les morphismes $s, t : R \to U$
et le morphisme $U' \to U$ sont universellement ouverts d'après
Morphismes, lemme \ref{morphisms-lemma-scheme-over-field-universally-open}.
Comme $R$ est non vide et $U$ est le spectre d'un corps, les morphismes
$s, t : R \to U$ sont surjectifs et plats. On conclut alors en utilisant
Morphismes d'espaces, lemmes
\ref{spaces-morphisms-lemma-base-change-surjective},
\ref{spaces-morphisms-lemma-composition-surjective},
\ref{spaces-morphisms-lemma-composition-open},
\ref{spaces-morphisms-lemma-base-change-affine},
\ref{spaces-morphisms-lemma-composition-affine},
\ref{spaces-morphisms-lemma-base-change-flat} et
\ref{spaces-morphisms-lemma-composition-flat}.
\end{proof}

\begin{lemma}
\label{lemma-groupoid-on-field-explain-points}
Dans la
situation \ref{situation-groupoid-on-field}.
Pour tout point $r \in |R|$, il existe
\begin{enumerate}
\item une extension de corps $k'/k$ avec $k'$ algébriquement clos,
\item un point $r' : \Spec(k') \to R'$, où
$(U', R', s', t', c')$ est la restriction de $(U, R, s, t, c)$
par $\Spec(k') \to \Spec(k)$,
\end{enumerate}
tels que
\begin{enumerate}
\item le point $r'$ s'envoie sur $r$ par le morphisme $R' \to R$, et
\item les applications
$s' \circ r', t' \circ r' : \Spec(k') \to \Spec(k')$
sont des automorphismes.
\end{enumerate}
\end{lemma}

\begin{proof}
Représentons $r$ par un morphisme $r : \Spec(K) \to R$ pour un certain
corps $K$. Pour démontrer le lemme, il faut trouver un corps
algébriquement clos $k'$ et un diagramme commutatif
$$
\xymatrix{
k' & k' \ar[l]^1 & \\
k' \ar[u]^\tau & K \ar[lu]^\sigma & k \ar[l]^-s \ar[lu]_i \\
& k \ar[lu]^i \ar[u]_t
}
$$
où $s, t : k \to K$ sont les homomorphismes de corps provenant de
$s \circ r$ et $t \circ r$. La démonstration de
Compléments sur les groupoïdes,
lemme \ref{more-groupoids-lemma-groupoid-on-field-explain-points},
explique comment construire un tel diagramme.
\end{proof}

\begin{lemma}
\label{lemma-groupoid-on-field-move-point}
Dans la
situation \ref{situation-groupoid-on-field}.
Si $r : \Spec(k) \to R$ est un morphisme tel que
$s \circ r, t \circ r$ soient des automorphismes de $\Spec(k)$, alors l'application
$$
R \longrightarrow R, \quad
x \longmapsto c(r, x)
$$
est un automorphisme $R \to R$ qui envoie $e$ sur $r$.
\end{lemma}

\begin{proof}
La démonstration est identique à celle de
Compléments sur les groupoïdes,
lemme \ref{more-groupoids-lemma-groupoid-on-field-move-point}.
\end{proof}

\begin{lemma}
\label{lemma-groupoid-on-field-geometrically-irreducible}
Dans la
situation \ref{situation-groupoid-on-field},
l'espace algébrique $R$ est géométriquement unibranche. Dans la
situation \ref{situation-group-over-field},
l'espace algébrique $G$ est géométriquement unibranche.
\end{lemma}

\begin{proof}
Soit $r \in |R|$. Il faut montrer que $R$ est géométriquement unibranche
en $r$. En combinant le
lemme \ref{lemma-restrict-groupoid-on-field}
avec
Descente sur les espaces, lemme \ref{spaces-descent-lemma-descend-unibranch},
on voit qu'il suffit de le démontrer lorsque $k$ est algébriquement clos
et que $r$ provient d'un morphisme $r : \Spec(k) \to R$ tel que
$s \circ r$ et $t \circ r$
soient des automorphismes de $\Spec(k)$. D'après le
lemme \ref{lemma-groupoid-on-field-move-point},
on se ramène au cas où $r = e$ est l'identité de $R$ et où $k$ est
algébriquement clos.

\medskip\noindent
Supposons que $r = e$ et que $k$ soit algébriquement clos. Soit
$A = \mathcal{O}_{R, e}$ l'anneau local étale de
$R$ en $e$, et soit
$C = \mathcal{O}_{R \times_{s, U, t} R, (e, e)}$
l'anneau local étale de $R \times_{s, U, t} R$ en $(e, e)$.
D'après Compléments d'algèbre, lemme
\ref{more-algebra-lemma-minimal-primes-tensor-strictly-henselian}
les idéaux premiers minimaux $\mathfrak q$ de $C$ correspondent $1$ à $1$
aux couples de premiers minimaux $\mathfrak p, \mathfrak p' \subset A$.
D'autre part, la loi de composition induit un homomorphisme plat d'anneaux
$$
\xymatrix{
A \ar[r]_{c^\sharp} & C & \mathfrak q \\
& A \otimes_{s^\sharp, k, t^\sharp} A \ar[u] &
\mathfrak p \otimes A + A \otimes \mathfrak p' \ar@{|}[u]
}
$$
Remarquons que $(c^\sharp)^{-1}(\mathfrak q)$ contient à la fois $\mathfrak p$ et
$\mathfrak p'$, puisque les diagrammes
$$
\xymatrix{
A \ar[r]_{c^\sharp} & C \\
A \otimes_{s^\sharp, k} k \ar[u] &
A \otimes_{s^\sharp, k, t^\sharp} A \ar[l]_{1 \otimes e^\sharp} \ar[u]
}
\quad\quad
\xymatrix{
A \ar[r]_{c^\sharp} & C \\
k \otimes_{k, t^\sharp} A \ar[u] &
A \otimes_{s^\sharp, k, t^\sharp} A \ar[l]_{e^\sharp \otimes 1} \ar[u]
}
$$
sont commutatifs d'après (\ref{equation-diagram}).
Comme $c^\sharp$ est plat (puisque $c$ est un morphisme plat d'après le
lemme \ref{lemma-groupoid-on-field-open-multiplication}),
on voit que $(c^\sharp)^{-1}(\mathfrak q)$ est un idéal premier minimal
de $A$. Ainsi $\mathfrak p = (c^\sharp)^{-1}(\mathfrak q) = \mathfrak p'$.
\end{proof}

\noindent
Dans le lemme suivant, nous utilisons la dimension des espaces algébriques
(en un point), telle qu'elle est définie dans
Propriétés des espaces, section \ref{spaces-properties-section-dimension}.
Nous utilisons aussi la dimension de l'anneau local définie dans
Propriétés des espaces, section
\ref{spaces-properties-section-dimension-local-ring},
ainsi que le degré de transcendance des points; voir
Morphismes d'espaces, section \ref{spaces-morphisms-section-relative-dimension}.

\begin{lemma}
\label{lemma-groupoid-on-field-locally-finite-type-dimension}
Dans la
situation \ref{situation-groupoid-on-field},
supposons $s, t$ localement de type fini.
Pour tout $r \in |R|$,
\begin{enumerate}
\item $\dim(R) = \dim_r(R)$,
\item le degré de transcendance de $r$ sur $\Spec(k)$
par $s$ est égal au degré de transcendance de $r$ sur $\Spec(k)$
par $t$, et
\item si le degré de transcendance mentionné en (2) est $0$, alors
$\dim(R) = \dim(\mathcal{O}_{R, \overline{r}})$.
\end{enumerate}
\end{lemma}

\begin{proof}
Soit $r \in |R|$. Notons $\text{trdeg}(r/_{\!\! s}k)$ le degré de
transcendance de $r$ sur $\Spec(k)$ par $s$. Choisissons un morphisme étale
$\varphi : V \to R$, où $V$ est un schéma et $v \in V$ s'envoie sur $r$.
Les définitions rappelées avant le lemme donnent
$$
\dim_r(R) = \dim_v(V) =
\dim(\mathcal{O}_{V, v}) + \text{trdeg}_{s(k)}(\kappa(v)) =
\dim(\mathcal{O}_{R, \overline{r}}) + \text{trdeg}(r/_{\!\! s}k)
$$
et de même pour $t$ (la deuxième égalité provenant de
Morphismes, lemme \ref{morphisms-lemma-dimension-fibre-at-a-point}).
On voit donc que $\text{trdeg}(r/_{\!\! s}k) = \text{trdeg}(r/_{\!\! t}k)$,
c'est-à-dire que (2) est vérifié.

\medskip\noindent
Soit $k'/k$ une extension de corps. Remarquons que la restriction $R'$
de $R$ à $\Spec(k')$ (voir le
lemme \ref{lemma-restrict-groupoid-on-field})
s'obtient à partir de $R$ par deux changements de base par des morphismes de corps. Ainsi,
Morphismes d'espaces,
lemme \ref{spaces-morphisms-lemma-dimension-fibre-after-base-change},
montre que la dimension de $R$ en un point reste inchangée par cette opération.
Pour démontrer (1), on peut donc supposer, d'après le
lemme \ref{lemma-groupoid-on-field-explain-points},
que $r$ est représenté par un morphisme $r : \Spec(k) \to R$ tel
que $s \circ r$ et $t \circ r$ soient tous deux des automorphismes de $\Spec(k)$.
Il existe alors un automorphisme $R \to R$ qui envoie $r$ sur $e$
(lemme \ref{lemma-groupoid-on-field-move-point}).
On voit donc que $\dim_r(R) = \dim_e(R)$ pour tout $r$. Par définition, cela
signifie que $\dim_r(R) = \dim(R)$.

\medskip\noindent
Le point (3) est une conséquence formelle des résultats obtenus dans la
discussion précédente.
\end{proof}

\begin{lemma}
\label{lemma-group-over-field-locally-finite-type-dimension}
Dans la
situation \ref{situation-group-over-field},
supposons $G$ localement de type fini.
Pour tout $g \in |G|$,
\begin{enumerate}
\item $\dim(G) = \dim_g(G)$,
\item si le degré de transcendance de $g$ sur $k$ est $0$, alors
$\dim(G) = \dim(\mathcal{O}_{G, \overline{g}})$.
\end{enumerate}
\end{lemma}

\begin{proof}
Cela résulte immédiatement du
lemme \ref{lemma-groupoid-on-field-locally-finite-type-dimension}
par (\ref{equation-groupoid-from-group}).
\end{proof}

\begin{lemma}
\label{lemma-groupoid-on-field-dimension-equal-stabilizer}
Dans la
situation \ref{situation-groupoid-on-field},
supposons $s, t$ localement de type fini.
Soit
$G = \Spec(k)
\times_{\Delta, \Spec(k) \times_B \Spec(k), t \times s} R$
l'espace algébrique en groupes stabilisateur.
Alors $\dim(R) = \dim(G)$.
\end{lemma}

\begin{proof}
Comme $G$ et $R$ sont équidimensionnels (voir les
lemmes \ref{lemma-groupoid-on-field-locally-finite-type-dimension} et
\ref{lemma-group-over-field-locally-finite-type-dimension})
il suffit de montrer que $\dim_e(R) = \dim_e(G)$. Soit $V$ un schéma affine,
soit $v \in V$, et soit $\varphi : V \to R$ un morphisme étale de schémas
tel que $\varphi(v) = e$. Remarquons que $V$ est noethérien, car
$s \circ \varphi$ est localement de type fini comme composée de morphismes
localement de type fini et parce que $V$ est quasi-compact (utiliser
Morphismes d'espaces, lemmes
\ref{spaces-morphisms-lemma-composition-finite-type},
\ref{spaces-morphisms-lemma-etale-locally-finite-presentation}, et
\ref{spaces-morphisms-lemma-finite-presentation-finite-type}
et
Morphismes, lemme \ref{morphisms-lemma-finite-type-noetherian}).
Ainsi, $V$ est localement connexe (voir
Propriétés, lemme \ref{properties-lemma-Noetherian-topology},
et
Topologie, lemme \ref{topology-lemma-locally-Noetherian-locally-connected}).
On peut donc remplacer $V$ par la composante connexe contenant $v$ (elle
reste affine, étant un sous-schéma ouvert et fermé de $V$).
Posons $T = V_{red}$, la réduction de $V$. Considérons les deux
morphismes $a, b : T \to \Spec(k)$ donnés par
$a = s \circ \varphi|_T$ et $b = t \circ \varphi|_T$. Remarquons que
$a, b$ induisent le même homomorphisme de corps $k \to \kappa(v)$, car $\varphi(v) = e$ !
Soit $k_a \subset \Gamma(T, \mathcal{O}_T)$ la clôture intégrale de
$a^\sharp(k) \subset \Gamma(T, \mathcal{O}_T)$. De même, soit
$k_b \subset \Gamma(T, \mathcal{O}_T)$ la clôture intégrale de
$b^\sharp(k) \subset \Gamma(T, \mathcal{O}_T)$. D'après
Variétés, proposition \ref{varieties-proposition-unique-base-field},
on a $k_a = k_b$. On obtient ainsi le diagramme commutatif suivant :
$$
\xymatrix{
k \ar[rd]^a \ar[rrrd] \\
& k_a = k_b \ar[r] & \Gamma(T, \mathcal{O}_T) \ar[r] & \kappa(v) \\
k \ar[ru]_b \ar[rrru]
}
$$
Comme on l'a vu ci-dessus, les longues flèches sont égales.
Puisque $k_a = k_b \to \kappa(v)$ est injectif, on en conclut que
les deux morphismes $a$ et $b$ sont égaux. Ainsi $T \to R$ se factorise par $G$.
Il s'ensuit que $R_{red} = G_{red}$ dans un voisinage ouvert de $e$,
ce qui implique bien que $\dim_e(R) = \dim_e(G)$.
\end{proof}






\section{Espaces algébriques en groupes sur les corps}
\label{section-group-algebraic-spaces-over-fields}

\noindent
Il existe un espace algébrique en groupes non séparé sur un corps,
à savoir $\mathbf{G}_a/\mathbf{Z}$ sur un corps de caractéristique zéro; voir
Exemples, section \ref{examples-section-non-separated-group-space}.
En fait, tout schéma en groupes sur un corps est séparé
(lemme \ref{lemma-group-scheme-over-field-separated});
ainsi, tout espace algébrique en groupes non séparé sur un corps
n'est pas représentable. D'autre part, un espace algébrique en
groupes sur un corps est séparé dès qu'il est décent; voir le
lemme \ref{lemma-group-scheme-over-field-separated}.
Dans cette section, nous allons montrer qu'un
espace algébrique en groupes séparé sur un corps
est représentable, c'est-à-dire qu'il est un schéma.

\begin{lemma}
\label{lemma-group-space-scheme-over-kbar}
Soit $k$ un corps muni d'une clôture algébrique $\overline{k}$.
Soit $G$ un espace algébrique en groupes sur $k$
qui est séparé\footnote{Il suffit de supposer $G$ décent,
par exemple localement séparé ou quasi-séparé d'après le
lemme \ref{lemma-group-scheme-over-field-separated}.}.
Alors $G_{\overline{k}}$ est un schéma.
\end{lemma}

\begin{proof}
D'après Espaces sur les corps, lemme
\ref{spaces-over-fields-lemma-when-scheme-after-base-change}
il suffit de montrer que $G_K$ est un schéma pour une certaine extension de corps
$K/k$. Notons $G_K' \subset G_K$ le lieu schématique
de $G_K$, comme dans Propriétés des espaces, lemme
\ref{spaces-properties-lemma-subscheme}.
D'après Propriétés des espaces, proposition
\ref{spaces-properties-proposition-locally-quasi-separated-open-dense-scheme}
on voit que $G_K' \subset G_K$ est un ouvert dense, en particulier non vide.
Choisissons un schéma $U$ et un morphisme étale surjectif $U \to G$. D'après
Variétés, lemme \ref{varieties-lemma-make-Jacobson},
si $K$ est un corps algébriquement clos de degré de transcendance assez
grand, alors $U_K$ est un schéma de Jacobson et
tout point fermé de $U_K$ est $K$-rationnel.
Ainsi $G_K'$ possède un point $K$-rationnel et il suffit
de montrer que tout point $K$-rationnel de $G_K$ appartient à $G_K'$.
Si $g \in G_K(K)$ est un point $K$-rationnel et $g' \in G_K'(K)$
un point $K$-rationnel du lieu schématique, alors on voit que
$g$ appartient à l'image de $G_K'$ par l'automorphisme
$$
G_K \longrightarrow G_K,\quad
h \longmapsto g(g')^{-1}h
$$
de $G_K$. Comme les automorphismes de l'espace algébrique $G_K$ préservent
$G_K'$, on conclut que $g \in G_K'$, comme voulu.
\end{proof}

\begin{lemma}
\label{lemma-group-space-scheme-locally-finite-type-over-k}
Soit $k$ un corps. Soit $G$ un espace algébrique en groupes sur $k$.
Si $G$ est séparé et localement de type fini sur $k$,
alors $G$ est un schéma.
\end{lemma}

\begin{proof}
Cela résulte du
lemme \ref{lemma-group-space-scheme-over-kbar},
de Groupoïdes, lemme \ref{groupoids-lemma-points-in-affine}, et
d'Espaces sur les corps, lemme
\ref{spaces-over-fields-lemma-scheme-over-algebraic-closure-enough-affines}.
\end{proof}

\begin{proposition}
\label{proposition-group-space-scheme-over-field}
Soit $k$ un corps. Soit $G$ un espace algébrique en groupes sur $k$.
Si $G$ est séparé, alors $G$ est un schéma.
\end{proposition}

\begin{proof}
Ce lemme généralise le
lemme \ref{lemma-group-space-scheme-locally-finite-type-over-k}
(qui couvre tous les cas utiles en pratique).
La démonstration est très semblable à celle d'
Espaces sur les corps, lemme
\ref{spaces-over-fields-lemma-scheme-over-algebraic-closure-enough-affines}
utilisé dans la démonstration du
lemme \ref{lemma-group-space-scheme-locally-finite-type-over-k};
nous invitons le lecteur à commencer par lire cette démonstration.

\medskip\noindent
D'après le lemme \ref{lemma-group-space-scheme-over-kbar}, le changement
de base $G_{\overline{k}}$ est un schéma.
Soit $K/k$ une extension purement transcendante de très grand
degré de transcendance. D'après Espaces sur les corps, lemme
\ref{spaces-over-fields-lemma-scheme-after-purely-transcendental-base-change}
il suffit de montrer que $G_K$ est un schéma.
Soit $K^{perf}$ la clôture parfaite de $K$. D'après
Espaces sur les corps, lemme
\ref{spaces-over-fields-lemma-scheme-after-purely-inseparable-base-change}
il suffit de montrer que $G_{K^{perf}}$ est un schéma.
Soit $K \subset K^{perf} \subset \overline{K}$, cette dernière étant la clôture algébrique
de $K$. On peut choisir un plongement
$\overline{k} \to \overline{K}$ au-dessus de $k$, de sorte que $G_{\overline{K}}$
soit le changement de base du schéma $G_{\overline{k}}$ par $\overline{k} \to \overline{K}$. D'après
Variétés, lemme \ref{varieties-lemma-make-Jacobson},
on voit que $G_{\overline{K}}$ est un schéma de Jacobson dont tous les
points fermés ont pour corps résiduel $\overline{K}$.

\medskip\noindent
Comme $G_{\overline{K}} \to G_{K^{perf}}$ est surjectif, il suffit de
montrer que l'image $g \in |G_{K^{perf}}|$ d'un point fermé arbitraire de
$G_{\overline{K}}$ appartient au lieu schématique de $G_K$.
En particulier, on peut représenter $g$ par un morphisme
$g : \Spec(L) \to G_{K^{perf}}$, où $L/K^{perf}$ est algébrique séparable
(on peut par exemple prendre $L = \overline{K}$). Le schéma
\begin{align*}
T & = \Spec(L) \times_{G_{K^{perf}}} G_{\overline{K}} \\
& =
\Spec(L) \times_{\Spec(K^{perf})} \Spec(\overline{K}) \\
& =
\Spec(L \otimes_{K^{perf}} \overline{K})
\end{align*}
est donc le spectre d'une $\overline{K}$-algèbre qui est une limite inductive
filtrante d'algèbres, produits finis de copies de $\overline{K}$.
Ainsi, d'après Groupoïdes, lemme \ref{groupoids-lemma-compact-set-in-affine},
on peut trouver un ouvert affine $W \subset G_{\overline{K}}$ contenant
l'image de $g_{\overline{K}} : T \to G_{\overline{K}}$.

\medskip\noindent
Choisissons un ouvert quasi-compact $V \subset G_{K^{perf}}$ contenant
l'image de $W$. D'après Espaces sur les corps,
lemme \ref{spaces-over-fields-lemma-when-scheme-after-base-change},
on voit que $V_{K'}$ est un schéma pour une certaine extension finie $K'/K^{perf}$.
En agrandissant $K'$, on peut supposer qu'il existe un ouvert affine
$U' \subset V_{K'} \subset G_{K'}$ dont le changement de base à $\overline{K}$
redonne $W$ (utiliser que $V_{\overline{K}}$ est la limite des schémas
$V_{K''}$ pour $K' \subset K'' \subset \overline{K}$ fini, ainsi que
Limites, lemmes \ref{limits-lemma-descend-opens} et
\ref{limits-lemma-limit-affine}). On peut supposer
que $K'/K^{perf}$ est une extension galoisienne (prendre la clôture normale,
Corps, lemme \ref{fields-lemma-normal-closure}, et utiliser
que $K^{perf}$ est parfait). Posons $H = \text{Gal}(K'/K^{perf})$.
Par construction, le sous-schéma fermé $H$-invariant
$\Spec(L) \times_{G_{K^{perf}}} G_{K'}$ est contenu dans $U'$.
D'après
Espaces sur les corps, lemmes
\ref{spaces-over-fields-lemma-base-change-by-Galois} et
\ref{spaces-over-fields-lemma-when-quotient-scheme-at-point}, on conclut.
\end{proof}





\section{Absence de courbes rationnelles sur les groupes}
\label{section-no-rational-curves}

\noindent
Dans cette section, nous montrons qu'il n'existe pas de morphisme non constant
de $\mathbf{P}^1$ vers un espace algébrique en groupes localement de type
fini sur un corps.

\begin{lemma}
\label{lemma-factor-through-over-open}
Soit $S$ un schéma. Soit $B$ un espace algébrique au-dessus de $S$.
Soient $f : X \to Y$ et $g : X \to Z$ des morphismes d'espaces
algébriques au-dessus de $B$. Supposons que
\begin{enumerate}
\item $Y \to B$ soit séparé,
\item $g$ soit surjectif, plat et localement de présentation finie,
\item il existe un ouvert schématiquement dense $V \subset Z$
tel que $f|_{g^{-1}(V)} : g^{-1}(V) \to Y$ se factorise par $V$.
\end{enumerate}
Alors $f$ se factorise par $g$.
\end{lemma}

\begin{proof}
Posons $R = X \times_Z X$. D'après (2), on a $Z = X/R$ comme faisceaux.
De plus, (2) implique que l'image inverse de $V$ dans $R$ est
schématiquement dense dans $R$ (Morphismes d'espaces, lemme
\ref{spaces-morphisms-lemma-flat-morphism-scheme-theoretically-dense-open}).
On voit alors que les deux composées
$R \to X \to Y$ sont égales d'après Morphismes d'espaces, lemme
\ref{spaces-morphisms-lemma-equality-of-morphisms}.
Le lemme en résulte.
\end{proof}

\begin{lemma}
\label{lemma-quotient-power-P1}
\begin{slogan}
Un morphisme d'un produit non vide de droites projectives sur un corps vers
un espace algébrique séparé de type fini sur un corps se factorise en un
morphisme fini après projection sur un produit de droites projectives.
\end{slogan}
Soit $k$ un corps. Soit $n \geq 1$ et soit $(\mathbf{P}^1_k)^n$
le produit de $n$ copies au-dessus de $\Spec(k)$. Soit
$f : (\mathbf{P}^1_k)^n \to Z$ un morphisme d'espaces algébriques sur $k$.
Si $Z$ est séparé et de type fini sur $k$, alors $f$ se factorise sous la forme
$$
(\mathbf{P}^1_k)^n \xrightarrow{projection}
(\mathbf{P}^1_k)^m \xrightarrow{fini} Z.
$$
\end{lemma}

\begin{proof}
On peut supposer $k$ algébriquement clos (détails omis); nous ne le faisons
que pour pouvoir raisonner avec des points rationnels, mais le lecteur peut
contourner cet artifice s'il le souhaite. Dans la preuve, les produits sont pris sur $k$.
L'espace algébrique en groupes des automorphismes de $(\mathbf{P}^1_k)^n$ contient
$G = (\text{GL}_{2, k})^n$. Si $C \subset (\mathbf{P}^1_k)^n$ est une
sous-variété fermée (donc en particulier irréductible sur $k$) envoyée
sur un point, on peut appliquer
Compléments sur les morphismes d'espaces,
lemme \ref{spaces-more-morphisms-lemma-flat-proper-family-cannot-collapse-fibre}
au morphisme
$$
G \times C \to G \times Z,\quad (g, c) \mapsto (g, f(g \cdot c))
$$
au-dessus de $G$. Ainsi, $g(C)$ est envoyé sur un point pour $g \in G(k)$
appartenant à un ouvert de Zariski $U \subset G$. Supposons que
$x = (x_1, \ldots, x_n)$, $y = (y_1, \ldots, y_n)$
soient des points à valeurs dans $k$ de $(\mathbf{P}^1_k)^n$. Soit
$I \subset \{1, \ldots, n\}$ l'ensemble des indices $i$
tels que $x_i = y_i$. Alors
$$
\{g(x) \mid g(y) = y,\ g \in U(k)\}
$$
est dense pour la topologie de Zariski dans la fibre de la projection
$\pi_I : (\mathbf{P}^1_k)^n \to \prod_{i \in I} \mathbf{P}^1_k$
(exercice). Ainsi, si $x, y \in C(k)$ sont distincts, on conclut
que $f$ envoie sur un seul point toute la fibre de $\pi_I$ contenant $x, y$.
De plus, l'orbite de $C$ sous $U(k)$ rencontre un ensemble ouvert de Zariski
de fibres de $\pi_I$. D'après le lemme \ref{lemma-factor-through-over-open},
le morphisme $f$ se factorise par $\pi_I$.
Après avoir répété ce procédé un nombre fini de fois, on atteint
le stade où toutes les fibres de $f$ au-dessus des points de $k$ sont finies.
Dans ce cas, $f$ est fini d'après
Compléments sur les morphismes d'espaces, lemme
\ref{spaces-more-morphisms-lemma-proper-finite-fibre-finite-in-neighbourhood}
et le fait que les points de $k$ sont denses dans $Z$
(Espaces sur les corps, lemme
\ref{spaces-over-fields-lemma-smooth-separable-closed-points-dense}).
\end{proof}

\begin{lemma}
\label{lemma-no-nonconstant-morphism-from-P1-to-group}
\begin{slogan}
Pas de courbes rationnelles complètes sur les groupes.
\end{slogan}
Soit $k$ un corps. Soit $G$ un espace algébrique en groupes séparé,
localement de type fini sur $k$. Il n'existe aucun
morphisme non constant $f : \mathbf{P}^1_k \to G$ au-dessus de $\Spec(k)$.
\end{lemma}

\begin{proof}
Supposons $f$ non constant. Considérons les morphismes
$$
\mathbf{P}^1_k \times_{\Spec(k)} \ldots \times_{\Spec(k)} \mathbf{P}^1_k
\longrightarrow G,
\quad (t_1, \ldots, t_n) \longmapsto f(t_1) \ldots f(t_n)
$$
où l'on utilise au membre de droite la multiplication du groupe.
D'après le lemme \ref{lemma-quotient-power-P1} et l'hypothèse que $f$
est non constant, ce morphisme est fini sur son image.
Ainsi $\dim(G) \geq n$ pour tout $n$, ce qui est impossible d'après le
lemme \ref{lemma-group-over-field-locally-finite-type-dimension}
et le fait que $G$ est localement de type fini sur $k$.
\end{proof}




\section{La partie finie d'un morphisme}
\label{section-finite}

\noindent
Soit $S$ un schéma.
Soit $f : X \to Y$ un morphisme d'espaces algébriques au-dessus de $S$.
Pour un espace algébrique ou un schéma $T$ au-dessus de $S$, considérons les couples
$(a, Z)$ où
\begin{equation}
\label{equation-finite-conditions}
\begin{matrix}
a : T \to Y\text{ est un morphisme au-dessus de }S, \\
Z \subset T \times_Y X\text{ est un sous-espace ouvert} \\
\text{tel que }\text{pr}_0|_Z : Z \to T\text{ soit fini.}
\end{matrix}
\end{equation}
Supposons que $h : T' \to T$ soit un morphisme d'espaces algébriques au-dessus de $S$
et que $(a, Z)$ soit un couple comme dans (\ref{equation-finite-conditions}) au-dessus de $T$. Posons
$a' = a \circ h$ et $Z' = (h \times \text{id}_X)^{-1}(Z) = T' \times_T Z$.
Alors $(a', Z')$ est un couple comme dans (\ref{equation-finite-conditions}) au-dessus de $T'$.
Cela résulte de ce que les morphismes finis sont préservés par changement de base; voir
Morphismes d'espaces, lemme \ref{spaces-morphisms-lemma-base-change-integral}.
On obtient ainsi un foncteur
\begin{equation}
\label{equation-finite}
\begin{matrix}
(X/Y)_{fin} : &
(\Sch/S)^{opp} &
\longrightarrow &
\textit{Ensembles} \\
& T & \longmapsto &
\{(a, Z)\text{ comme ci-dessus}\}
\end{matrix}
\end{equation}
Pour les applications, nous nous intéressons surtout à ce foncteur $(X/Y)_{fin}$
lorsque $f$ est séparé et localement de type fini. Pour se faire une idée
de la construction, consulter la
remarque \ref{remark-finite-quasi-finite-separated-morphism-schemes}.

\begin{lemma}
\label{lemma-finite-sheaf}
Soit $S$ un schéma.
Soit $f : X \to Y$ un morphisme d'espaces algébriques au-dessus de $S$.
Alors
\begin{enumerate}
\item Le préfaisceau $(X/Y)_{fin}$ satisfait la condition de faisceau pour
la topologie fppf.
\item Si $T$ est un espace algébrique au-dessus de $S$, il existe une
bijection canonique
$$
\Mor_{\Sh((\Sch/S)_{fppf})}(T, (X/Y)_{fin})
=
\{(a, Z)\text{ satisfaisant \ref{equation-finite-conditions}}\}
$$
\end{enumerate}
\end{lemma}

\begin{proof}
Soit $T$ un espace algébrique au-dessus de $S$.
Soit $\{T_i \to T\}$ un recouvrement fppf (par des espaces algébriques).
Soient $s_i = (a_i, Z_i)$ des couples au-dessus de $T_i$
satisfaisant \ref{equation-finite-conditions}
tels que $s_i|_{T_i \times_T T_j} = s_j|_{T_i \times_T T_j}$.
Premièrement, cela implique notamment que $a_i$ et $a_j$ définissent le même
morphisme $T_i \times_T T_j \to Y$. D'après
Descente sur les espaces,
lemme \ref{spaces-descent-lemma-fpqc-universal-effective-epimorphisms},
on en déduit qu'il existe un unique morphisme $a : T \to Y$
tel que $a_i$ soit la composée $T_i \to T \to Y$.
Deuxièmement, cela implique que les $Z_i \subset T_i \times_Y X$ sont des sous-espaces
ouverts dont les images inverses dans $(T_i \times_T T_j) \times_Y X$ sont égales.
Comme $\{T_i \times_Y X \to T \times_Y X\}$ est un recouvrement fppf,
on en déduit qu'il existe un unique sous-espace ouvert $Z \subset T \times_Y X$
dont la restriction à $T_i$ est $Z_i$; voir
Descente sur les espaces, lemme \ref{spaces-descent-lemma-open-fpqc-covering}.
Nous affirmons que la projection $Z \to T$ est finie.
Cela résulte de ce que la finitude est locale pour la topologie fpqc; voir
Descente sur les espaces, lemme \ref{spaces-descent-lemma-descending-property-finite}.

\medskip\noindent
Remarquons que le résultat du paragraphe précédent implique notamment (1).

\medskip\noindent
Soit $T$ un espace algébrique au-dessus de $S$. Pour démontrer (2), nous allons
construire des applications réciproques entre les ensembles affichés. Dans la
suite, le mot « couple » désigne un couple satisfaisant les
conditions \ref{equation-finite-conditions}.

\medskip\noindent
Soit $v : T \to (X/Y)_{fin}$ une transformation naturelle.
Choisissons un schéma $U$ et un morphisme étale surjectif $p : U \to T$.
Alors $v(p) \in (X/Y)_{fin}(U)$ correspond à un couple $(a_U, Z_U)$
au-dessus de $U$. Soit $R = U \times_T U$, de projections $t, s : R \to U$.
Comme $v$ est une transformation de foncteurs, les images inverses de
$(a_U, Z_U)$ par $s$ et $t$ sont égales. Puisque $\{U \to T\}$ est un
recouvrement fppf, on peut donc appliquer le résultat du premier paragraphe et
en déduire qu'il existe un unique couple $(a, Z)$ au-dessus de $T$.

\medskip\noindent
Réciproquement, soit $(a, Z)$ un couple au-dessus de $T$.
Soient $U \to T$, $R = U \times_T U$ et $t, s : R \to U$ comme
ci-dessus. La restriction $(a, Z)|_U$ donne alors naissance à une
transformation de foncteurs $v : h_U \to (X/Y)_{fin}$ par le
lemme de Yoneda
(Catégories, lemme \ref{categories-lemma-yoneda}).
Comme les deux images inverses $s^*(a, Z)|_U$ et $t^*(a, Z)|_U$
sont égales, on voit que $v$ coégalise les deux applications
$h_t, h_s : h_R \to h_U$. Puisque $T = U/R$ est le faisceau quotient fppf d'après
Espaces, lemme \ref{spaces-lemma-space-presentation},
et que $(X/Y)_{fin}$ est un faisceau fppf d'après (1), on conclut
que $v$ se factorise par une application $T \to (X/Y)_{fin}$.

\medskip\noindent
Nous omettons la vérification que les deux constructions ci-dessus sont
réciproques.
\end{proof}

\begin{lemma}
\label{lemma-finite-open}
Soit $S$ un schéma. Considérons un diagramme commutatif
$$
\xymatrix{
X' \ar[rr]_j \ar[rd] & & X \ar[ld] \\
& Y
}
$$
d'espaces algébriques au-dessus de $S$. Si $j$ est une immersion ouverte, alors
il existe une application injective canonique de faisceaux
$j : (X'/Y)_{fin} \to (X/Y)_{fin}$.
\end{lemma}

\begin{proof}
Si $(a, Z)$ est un couple au-dessus de $T$ pour $X'/Y$, alors
$(a, j(Z))$ est un couple au-dessus de $T$ pour $X/Y$.
\end{proof}

\begin{lemma}
\label{lemma-finite-lives-on-locally-quasi-finite-part}
Soit $S$ un schéma.
Soit $f : X \to Y$ un morphisme d'espaces algébriques au-dessus de $S$ qui est
localement de type fini.
Soit $X' \subset X$ le plus grand sous-espace ouvert sur lequel $f$ est
localement quasi-fini; voir
Morphismes d'espaces,
lemme \ref{spaces-morphisms-lemma-locally-finite-type-quasi-finite-part}.
Alors $(X/Y)_{fin} = (X'/Y)_{fin}$.
\end{lemma}

\begin{proof}
Le lemme \ref{lemma-finite-open}
fournit une application injective $(X'/Y)_{fin} \to (X/Y)_{fin}$.
Morphismes d'espaces,
lemme \ref{spaces-morphisms-lemma-locally-finite-type-quasi-finite-part},
assure que la formation de $X'$ commute au changement de base.
Il suffit donc de montrer que, si
$Z \subset X$ est un sous-espace ouvert tel que $f|_Z : Z \to Y$ soit fini,
alors $Z \subset X'$. Cela est vrai puisqu'un morphisme fini
est localement quasi-fini; voir
Morphismes d'espaces, lemme \ref{spaces-morphisms-lemma-finite-quasi-finite}.
\end{proof}

\begin{lemma}
\label{lemma-finite-separated}
Soit $S$ un schéma.
Soit $f : X \to Y$ un morphisme d'espaces algébriques au-dessus de $S$.
Soit $T$ un espace algébrique au-dessus de $S$, et soit $(a, Z)$
un couple comme dans \ref{equation-finite-conditions}.
Si $f$ est séparé, alors $Z$ est fermé dans $T \times_Y X$.
\end{lemma}

\begin{proof}
Un morphisme fini d'espaces algébriques est universellement fermé d'après
Morphismes d'espaces, lemme \ref{spaces-morphisms-lemma-finite-proper}.
Comme $f$ est séparé, le morphisme $T \times_Y X \to T$ l'est aussi; voir
Morphismes d'espaces, lemme \ref{spaces-morphisms-lemma-base-change-separated}.
Le caractère fermé de $Z$ résulte donc de
Morphismes d'espaces,
lemme \ref{spaces-morphisms-lemma-universally-closed-permanence}.
\end{proof}

\begin{remark}
\label{remark-finite-monoid}
Soit $f : X \to Y$ un morphisme séparé d'espaces algébriques.
Le faisceau $(X/Y)_{fin}$ est muni d'une application naturelle
$(X/Y)_{fin} \to Y$ qui envoie le couple $(a, Z) \in (X/Y)_{fin}(T)$
sur l'élément $a \in Y(T)$. On peut utiliser le
lemme \ref{lemma-finite-separated}
pour définir les opérations
$$
\star_i : (X/Y)_{fin} \times_Y (X/Y)_{fin} \longrightarrow (X/Y)_{fin}
$$
par les règles
\begin{align*}
\star_1 : ((a, Z_1), (a, Z_2)) & \longmapsto (a, Z_1 \cup Z_2) \\
\star_2 : ((a, Z_1), (a, Z_2)) & \longmapsto (a, Z_1 \cap Z_2) \\
\star_3 : ((a, Z_1), (a, Z_2)) & \longmapsto (a, Z_1 \setminus Z_2) \\
\star_4 : ((a, Z_1), (a, Z_2)) & \longmapsto (a, Z_2 \setminus Z_1).
\end{align*}
Cette construction fonctionne parce que $Z_1 \cap Z_2$ est à la fois ouvert et fermé
dans $Z_1$ et $Z_2$ (ce qui implique aussi que $Z_1 \cup Z_2$ est
la réunion disjointe des trois autres morceaux).
On peut donc considérer $(X/Y)_{fin}$ comme une $\mathbf{F}_2$-algèbre
(sans unité) au-dessus de $Y$, dont la multiplication est donnée par
$ss' = \star_2(s, s')$ et l'addition par
$$
s + s' = \star_1(\star_3(s, s'), \star_4(s, s'))
$$
ce qui revient à prendre la différence symétrique.
Remarquons que, dans ce faisceau d'algèbres, $0 = (1_Y, \emptyset)$
et que l'on a bien $s + s = 0$ pour toute section locale $s$.
Si $f : X \to Y$ est fini, cette algèbre possède une unité, à savoir
$1 = (1_Y, X)$, et $\star_3(s, s') = s(1 + s')$, tandis que
$\star_4(s, s') = (1 + s)s'$.
\end{remark}

\begin{remark}
\label{remark-finite-quasi-finite-separated-morphism-schemes}
Soit $f : X \to Y$ un morphisme de schémas séparé et localement
quasi-fini. Dans ce cas, le faisceau $(X/Y)_{fin}$
est étroitement lié au faisceau $f_!\mathbf{F}_2$
(insérer ici une référence future) sur $Y_\etale$.
En effet, si $V \to Y$ est étale et $s \in \Gamma(V, f_!\mathbf{F}_2)$,
alors $s \in \Gamma(V \times_Y X, \mathbf{F}_2)$ est une section
de support propre $Z = \text{Supp}(s)$ sur $V$. Comme $f$ est
aussi localement quasi-fini, la projection $Z \to V$ est en fait
finie. Comme le support d'une section d'un faisceau abélien constant est ouvert,
on voit que le couple $(V \to Y, \text{Supp}(s))$ satisfait
\ref{equation-finite-conditions}.
En fait, $f_!\mathbf{F}_2 \cong (X/Y)_{fin}|_{Y_\etale}$
dans ce cas, ce qui explique aussi la structure de $\mathbf{F}_2$-algèbre
introduite dans la remarque \ref{remark-finite-monoid}.
\end{remark}

\begin{lemma}
\label{lemma-finite-diagonal}
Soit $S$ un schéma.
Soit $f : X \to Y$ un morphisme d'espaces algébriques au-dessus de $S$.
La diagonale de $(X/Y)_{fin} \to Y$
$$
(X/Y)_{fin} \longrightarrow (X/Y)_{fin} \times_Y (X/Y)_{fin}
$$
est représentable (par des schémas) et est une immersion ouverte; la diagonale
« absolue »
$$
(X/Y)_{fin} \longrightarrow (X/Y)_{fin} \times (X/Y)_{fin}
$$
est représentable (par des schémas).
\end{lemma}

\begin{proof}
La seconde assertion résulte de la première, car la diagonale absolue
est la composée de la diagonale relative et d'un changement de base
de la diagonale de $Y$ (qui est représentable par des schémas); voir
Espaces, section \ref{spaces-section-representable}.
Pour démontrer la première assertion, il faut établir ce qui suit.
Étant donnés un schéma $T$ et deux couples $(a, Z_1)$ et $(a, Z_2)$ au-dessus de $T$
ayant la même première composante et
satisfaisant \ref{equation-finite-conditions},
il existe un sous-schéma ouvert $V \subset T$ possédant la propriété suivante :
pour tout morphisme de schémas $h : T' \to T$, on a
$$
h(T') \subset V \Leftrightarrow
\Big(T' \times_T Z_1 = T' \times_T Z_2
\text{ comme sous-espaces de }T' \times_Y X\Big)
$$
Construisons $V$. Remarquons que $Z_1 \cap Z_2$ est ouvert dans $Z_1$
et dans $Z_2$. Comme $\text{pr}_0|_{Z_i} : Z_i \to T$ est fini,
donc propre (voir
Morphismes d'espaces, lemme \ref{spaces-morphisms-lemma-finite-proper}),
on voit que
$$
E =
\text{pr}_0|_{Z_1}\left(Z_1 \setminus Z_1 \cap Z_2)\right)
\cup
\text{pr}_0|_{Z_2}\left(Z_2 \setminus Z_1 \cap Z_2)\right)
$$
est fermé dans $T$. Il est alors clair que $V = T \setminus E$ convient.
\end{proof}

\begin{lemma}
\label{lemma-finite-criterion-etale}
Soit $S$ un schéma.
Soit $f : X \to Y$ un morphisme d'espaces algébriques au-dessus de $S$.
Supposons que $U$ soit un schéma, que $U \to Y$ soit un morphisme étale et que
$Z \subset U \times_Y X$ soit un sous-espace ouvert fini sur $U$.
Alors le morphisme induit $U \to (X/Y)_{fin}$ est étale.
\end{lemma}

\begin{proof}
Cela résulte formellement de la description de la diagonale dans le
lemme \ref{lemma-finite-diagonal},
mais nous en donnons les détails, car il s'agit d'une étape importante dans le
développement de la théorie. Il faut vérifier que, pour tout schéma $T$ au-dessus de $S$ et tout morphisme
$T \to (X/Y)_{fin}$, la projection
$$
T \times_{(X/Y)_{fin}} U \longrightarrow T
$$
est étale. Remarquons que
$$
T \times_{(X/Y)_{fin}} U
=
(X/Y)_{fin} \times_{((X/Y)_{fin} \times_Y (X/Y)_{fin})} (T \times_Y U)
$$
En appliquant le résultat du
lemme \ref{lemma-finite-diagonal},
on voit que $T \times_{(X/Y)_{fin}} U$ est représenté par un sous-schéma ouvert de
$T \times_Y U$. Comme la projection $T \times_Y U \to T$ est étale d'après
Morphismes d'espaces, lemme \ref{spaces-morphisms-lemma-base-change-etale},
on conclut.
\end{proof}

\begin{lemma}
\label{lemma-finite-pullback}
Soit $S$ un schéma.
Soit
$$
\xymatrix{
X' \ar[d] \ar[r] & X \ar[d] \\
Y' \ar[r] & Y
}
$$
un carré de produits fibrés d'espaces algébriques au-dessus de $S$. Alors
$$
\xymatrix{
(X'/Y')_{fin} \ar[d] \ar[r] & (X/Y)_{fin} \ar[d] \\
Y' \ar[r] & Y
}
$$
est un carré de produits fibrés de faisceaux sur $(\Sch/S)_{fppf}$.
\end{lemma}

\begin{proof}
Il résulte immédiatement des définitions que
le faisceau $(X'/Y')_{fin}$ est égal au faisceau
$Y' \times_Y (X/Y)_{fin}$.
\end{proof}

\begin{lemma}
\label{lemma-finite-surjective-etale-cover}
Soit $S$ un schéma.
Soit $f : X \to Y$ un morphisme d'espaces algébriques au-dessus de $S$.
Si $f$ est séparé et localement quasi-fini, il existe un
schéma $U$ étale sur $Y$ et un morphisme étale surjectif
$U \to (X/Y)_{fin}$ au-dessus de $Y$.
\end{lemma}

\begin{proof}
Remarquons que l'assertion a un sens d'après le résultat du
lemme \ref{lemma-finite-diagonal}
sur la diagonale de $(X/Y)_{fin}$; voir
Espaces, lemme \ref{spaces-lemma-representable-diagonal}.
Soit $V$ un schéma et soit $V \to Y$ un morphisme étale surjectif. D'après le
lemme \ref{lemma-finite-pullback},
le morphisme $(V \times_Y X/V)_{fin} \to (X/Y)_{fin}$ est
un changement de base de l'application $V \to Y$ et est donc surjectif et étale; voir
Espaces,
lemme \ref{spaces-lemma-base-change-representable-transformations-property}.
Il suffit donc de démontrer le lemme pour $(V \times_Y X/V)_{fin}$.
(On utilise ici implicitement que la composée de transformations de foncteurs
représentables, surjectives et étales est encore représentable, surjective et
étale; voir
Espaces, lemmes \ref{spaces-lemma-composition-representable-transformations} et
\ref{spaces-lemma-composition-representable-transformations-property}, ainsi que
Morphismes, lemmes \ref{morphisms-lemma-composition-surjective} et
\ref{morphisms-lemma-composition-etale}.)
Remarquons que les propriétés d'être séparé et localement quasi-fini
sont préservées par changement de base; voir
Morphismes d'espaces,
lemmes \ref{spaces-morphisms-lemma-base-change-separated} et
\ref{spaces-morphisms-lemma-base-change-quasi-finite}.
Ainsi $V \times_Y X \to V$ est lui aussi séparé et localement quasi-fini,
et, d'après
Morphismes d'espaces, proposition
\ref{spaces-morphisms-proposition-locally-quasi-finite-separated-over-scheme}
on voit que $V \times_Y X$ est également un schéma.
On peut donc supposer que $f : X \to Y$ est un morphisme de schémas
séparé et localement quasi-fini.

\medskip\noindent
Choisissons un point $y \in Y$. Choisissons des points $x_1, \ldots, x_n \in X$
au-dessus de $y$. Choisissons un voisinage étale $a : (U, u) \to (Y, y)$ et une
décomposition
$$
U \times_S X =
W \amalg
\ \coprod\nolimits_{i = 1, \ldots, n}
\ \coprod\nolimits_{j = 1, \ldots, m_j}
V_{i, j}
$$
comme dans
Compléments sur les morphismes, lemme
\ref{more-morphisms-lemma-etale-splits-off-quasi-finite-part-technical-variant}.
Choisissons un sous-ensemble quelconque
$$
I \subset \{(i, j) \mid 1 \leq i \leq n, \ 1 \leq j \leq m_i\}.
$$
Ces choix donnent un couple $(a, Z)$ avec
$Z = \bigcup_{(i, j) \in I} V_{i, j}$
qui satisfait les conditions \ref{equation-finite-conditions}. Autrement dit,
on obtient un morphisme $U \to (X/Y)_{fin}$. La construction de ce morphisme
dépend de tous les choix effectués ci-dessus, de sorte qu'il faudrait écrire
$$
U(y, n, x_1, \ldots, x_n, a, I) \longrightarrow (X/Y)_{fin}
$$
Ce morphisme est étale d'après le lemme \ref{lemma-finite-criterion-etale}.

\medskip\noindent
Affirmation : leur réunion disjointe est surjective sur $(X/Y)_{fin}$.
Il est clair que, si l'affirmation est vraie, le lemme en résulte.

\medskip\noindent
Pour établir la surjectivité, il faut montrer ce qui suit (voir
Espaces, remarque \ref{spaces-remark-warning}) : étant donnés un schéma $T$ au-dessus de
$S$, un point $t \in T$ et une application $T \to (X/Y)_{fin}$, on peut trouver une donnée
$(y, n, x_1, \ldots, x_n, a, I)$ comme ci-dessus telle que
$t$ appartienne à l'image de la projection
$$
U(y, n, x_1, \ldots, x_n, a, I) \times_{(X/Y)_{fin}} T \longrightarrow T.
$$
Pour le démontrer, on peut clairement remplacer $T$ par
$\Spec(\overline{\kappa(t)})$ et
$T \to (X/Y)_{fin}$ par la composée
$\Spec(\overline{\kappa(t)}) \to T \to (X/Y)_{fin}$.
Autrement dit, on peut supposer que $T$ est
le spectre d'un corps algébriquement clos.

\medskip\noindent
Soit $T = \Spec(k)$ le spectre d'un corps
algébriquement clos $k$. Le morphisme $T \to (X/Y)_{fin}$ est donné par un couple
$(T \to Y, Z)$ satisfaisant les conditions \ref{equation-finite-conditions}.
Voici un diagramme :
$$
\xymatrix{
& Z \ar[d] \ar[r] & X \ar[d] \\
\Spec(k) \ar@{=}[r] & T \ar[r] & Y
}
$$
Soit $y \in Y$ le point image de $T \to Y$.
Comme $Z$ est fini sur $k$, il possède un nombre fini de points.
Il existe donc des points $x_1, \ldots, x_n \in X$ en nombre fini
tels que l'image de $Z$ dans $X$ soit contenue dans $\{x_1, \ldots, x_n\}$.
Choisissons $a : (U, u) \to (Y, y)$ adapté à $y$ et $x_1, \ldots, x_n$ comme
ci-dessus; on obtient le diagramme
$$
\xymatrix{
W \amalg
\ \coprod\nolimits_{i = 1, \ldots, n}
\ \coprod\nolimits_{j = 1, \ldots, m_j}
V_{i, j} \ar[d] \ar[r] &
X \ar[d] \\
U \ar[r] & Y.
}
$$
Comme $k$ est algébriquement clos et que
$\kappa(y) \subset \kappa(u)$ est une extension finie séparable,
on peut factoriser le morphisme
$T = \Spec(k) \to Y$ par le morphisme
$u = \Spec(\kappa(u)) \to \Spec(\kappa(y)) = y \subset Y$.
Avec ce choix, on obtient le diagramme commutatif :
$$
\xymatrix{
Z \ar[d] \ar[r] &
W \amalg
\ \coprod\nolimits_{i = 1, \ldots, n}
\ \coprod\nolimits_{j = 1, \ldots, m_j}
V_{i, j} \ar[d] \ar[r] &
X \ar[d] \\
\Spec(k) \ar[r] &
U \ar[r] & Y
}
$$
Nous savons que l'image de la flèche supérieure gauche est contenue dans
$\coprod V_{i, j}$. Rappelons aussi que $Z$ est un sous-schéma ouvert
de $\Spec(k) \times_Y X$ par définition de $(X/Y)_{fin}$
et que le carré de droite est un carré de produits fibrés.
On voit donc que
$$
Z \subset
\coprod\nolimits_{i = 1, \ldots, n}\ \coprod\nolimits_{j = 1, \ldots, m_j}
\Spec(k) \times_U V_{i, j}
$$
est un sous-schéma ouvert. Par construction (voir
Compléments sur les morphismes, lemme
\ref{more-morphisms-lemma-etale-splits-off-quasi-finite-part-technical-variant})
chaque $V_{i, j}$ possède un unique point $v_{i, j}$ au-dessus de $u$
dont l'extension de corps résiduels
$\kappa(v_{i, j})/\kappa(u)$ est purement inséparable. Ainsi, chaque
schéma $\Spec(k) \times_U V_{i, j}$ possède exactement un
point. On voit donc que
$$
Z = \coprod\nolimits_{(i, j) \in I} \Spec(k) \times_U V_{i, j}
$$
pour un unique sous-ensemble
$I \subset \{(i, j) \mid 1 \leq i \leq n, \ 1 \leq j \leq m_i\}$.
En explicitant les définitions, cela montre que
$$
U(y, n, x_1, \ldots, x_n, a, I) \times_{(X/Y)_{fin}} T
$$
avec le $I$ trouvé ci-dessus est non vide, comme voulu.
\end{proof}

\begin{proposition}
\label{proposition-finite-algebraic-space}
Soit $S$ un schéma.
Soit $f : X \to Y$ un morphisme d'espaces algébriques au-dessus de $S$ qui
est séparé et localement de type fini. Alors $(X/Y)_{fin}$
est un espace algébrique. De plus, le morphisme
$(X/Y)_{fin} \to Y$ est étale.
\end{proposition}

\begin{proof}
D'après le
lemme \ref{lemma-finite-lives-on-locally-quasi-finite-part},
on peut remplacer $X$ par le sous-schéma ouvert localement quasi-fini
sur $Y$. On peut donc supposer $f$ séparé et localement quasi-fini.
Vérifions les trois conditions de
Espaces, définition \ref{spaces-definition-algebraic-space}.
La condition (1) résulte du
lemme \ref{lemma-finite-sheaf}.
La condition (2) résulte du
lemme \ref{lemma-finite-diagonal}.
Enfin, la condition (3) résulte du
lemme \ref{lemma-finite-surjective-etale-cover}.
Ainsi $(X/Y)_{fin}$ est un espace algébrique.
De plus, ce lemme montre qu'il existe un
diagramme commutatif
$$
\xymatrix{
U \ar[rr] \ar[rd] & & (X/Y)_{fin} \ar[ld] \\
& Y
}
$$
dont la flèche horizontale est surjective et étale et la flèche sud-est
est étale. D'après
Propriétés des espaces, lemme \ref{spaces-properties-lemma-etale-local},
cela implique que la flèche sud-ouest est elle aussi étale.
\end{proof}

\begin{remark}
\label{remark-warning}
L'hypothèse que $f$ soit séparé ne peut être supprimée dans la
proposition \ref{proposition-finite-algebraic-space}.
Un exemple consiste à prendre pour $X$ la droite affine à origine dédoublée; voir
Schémas, exemple \ref{schemes-example-affine-space-zero-doubled},
pour $Y = \mathbf{A}^1_k$ la droite affine, et pour $X \to Y$ l'application évidente.
Rappelons qu'au-dessus de $0 \in Y$, il y a deux points $0_1$ et $0_2$
dans $X$. Ainsi, $(X/Y)_{fin}$ possède quatre points au-dessus de $0$, à savoir
$\emptyset, \{0_1\}, \{0_2\}, \{0_1, 0_2\}$.
De ces quatre points, seuls trois peuvent se relever en un sous-schéma ouvert
de $U \times_Y X$ fini sur $U$ lorsque $U \to Y$ est étale,
à savoir $\emptyset, \{0_1\}, \{0_2\}$. Cela montre que $(X/Y)_{fin}$,
s'il est représentable par un espace algébrique, n'est pas étale sur $Y$.
Des arguments analogues montrent que $(X/Y)_{fin}$ n'est réellement pas un espace
algébrique. Les détails sont omis.
\end{remark}

\begin{remark}
\label{remark-not-scheme}
Soit $Y = \mathbf{A}^1_{\mathbf{R}}$ la droite affine sur le corps des nombres
réels, et soit $X = \Spec(\mathbf{C})$ envoyé sur le point
$\mathbf{R}$-rationnel $0$ de $Y$. Dans ce cas, le morphisme
$f : X \to Y$ est fini, mais $(X/Y)_{fin}$ n'est pas
un schéma. En effet, on peut montrer que, dans ce cas, l'espace
algébrique $(X/Y)_{fin}$ est isomorphe à l'espace algébrique d'
Espaces, exemple \ref{spaces-example-non-representable-descent},
associé à l'extension $\mathbf{R} \subset \mathbf{C}$.
Il est donc réellement nécessaire de sortir de la catégorie des schémas
pour représenter le faisceau $(X/Y)_{fin}$, même lorsque $f$
est un morphisme fini.
\end{remark}

\begin{lemma}
\label{lemma-finite-separated-flat-locally-finite-presentation}
Soit $S$ un schéma.
Soit $f : X \to Y$ un morphisme d'espaces algébriques au-dessus de $S$ qui
est séparé, plat et localement de présentation finie.
Dans ce cas,
\begin{enumerate}
\item $(X/Y)_{fin} \to Y$ est séparé, représentable et étale, et
\item si $Y$ est un schéma, alors $(X/Y)_{fin}$ est (représenté par) un schéma.
\end{enumerate}
\end{lemma}

\begin{proof}
Comme $f$ est en particulier séparé et localement de type fini (voir
Morphismes d'espaces,
lemme \ref{spaces-morphisms-lemma-finite-presentation-finite-type}),
on voit que $(X/Y)_{fin}$ est un espace algébrique d'après la
proposition \ref{proposition-finite-algebraic-space}.
Pour montrer que $(X/Y)_{fin} \to Y$ est séparé, il faut établir
ce qui suit : étant donnés un schéma $T$ et deux couples $(a, Z_1)$ et $(a, Z_2)$
au-dessus de $T$
ayant la même première composante et satisfaisant \ref{equation-finite-conditions},
il existe un sous-schéma fermé $V \subset T$ possédant la propriété suivante :
pour tout morphisme de schémas $h : T' \to T$, on a
$$
h \text{ se factorise par } V \Leftrightarrow
\Big(T' \times_T Z_1 = T' \times_T Z_2
\text{ comme sous-espaces de }T' \times_Y X\Big)
$$
Dans la démonstration du
lemme \ref{lemma-finite-diagonal},
on a vu que $V = T \setminus E$ est un sous-schéma ouvert de $T$
de complémentaire fermé
$$
E =
\text{pr}_0|_{Z_1}\left(Z_1 \setminus Z_1 \cap Z_2)\right)
\cup
\text{pr}_0|_{Z_2}\left(Z_2 \setminus Z_1 \cap Z_2)\right).
$$
Il suffit donc de montrer que $E$ est aussi ouvert. D'après le
lemme \ref{lemma-finite-separated},
on voit que $Z_1$ et $Z_2$ sont fermés dans $T \times_Y X$. Ainsi,
$Z_1 \setminus Z_1 \cap Z_2$ est ouvert dans $Z_1$. Comme $f$ est plat et
localement de présentation finie, $\text{pr}_0|_{Z_1}$ l'est aussi.
Cela est vrai parce que $Z_1$ est un sous-espace ouvert du changement de base
$T \times_Y X$, et d'après
Morphismes d'espaces,
lemmes \ref{spaces-morphisms-lemma-base-change-finite-presentation} et
lemme \ref{spaces-morphisms-lemma-base-change-flat}.
Ainsi $\text{pr}_0|_{Z_1}$ est ouvert; voir
Morphismes d'espaces, lemme \ref{spaces-morphisms-lemma-fppf-open}.
Par conséquent, $\text{pr}_0|_{Z_1}\left(Z_1 \setminus Z_1 \cap Z_2)\right)$ est
ouvert et il s'ensuit que $E$ est ouvert, comme voulu.

\medskip\noindent
Nous avons déjà vu que $(X/Y)_{fin} \to Y$ est étale; voir la
proposition \ref{proposition-finite-algebraic-space}.
Nous savons donc qu'il est localement quasi-fini (voir
Morphismes d'espaces,
lemme \ref{spaces-morphisms-lemma-etale-locally-quasi-finite})
et séparé, donc représentable d'après
Morphismes d'espaces, lemme
\ref{spaces-morphisms-lemma-locally-quasi-finite-separated-representable}.
La dernière assertion est claire (on peut, si l'on veut, utiliser
Morphismes d'espaces, proposition
\ref{spaces-morphisms-proposition-locally-quasi-finite-separated-over-scheme}).
\end{proof}

\noindent
Variante : soit $S$ un schéma.
Soit $f : X \to Y$ un morphisme d'espaces algébriques au-dessus de $S$.
Soit $\sigma : Y \to X$ une section de $f$.
Pour un espace algébrique ou un schéma $T$ au-dessus de $S$, considérons les couples
$(a, Z)$ où
\begin{equation}
\label{equation-finite-conditions-variant}
\begin{matrix}
a : T \to Y\text{ est un morphisme au-dessus de }S, \\
Z \subset T \times_Y X\text{ est un sous-espace ouvert} \\
\text{tel que }\text{pr}_0|_Z : Z \to T\text{ soit fini et} \\
(1_T, \sigma \circ a) : T \to T \times_Y X\text{ se factorise par }Z.
\end{matrix}
\end{equation}
Nous noterons $(X/Y, \sigma)_{fin}$ le sous-foncteur de $(X/Y)_{fin}$
qui paramètre ces couples.

\begin{lemma}
\label{lemma-finite-plus-section}
Soit $S$ un schéma.
Soit $f : X \to Y$ un morphisme d'espaces algébriques au-dessus de $S$.
Soit $\sigma : Y \to X$ une section de $f$. Considérons la
transformation de foncteurs
$$
t : (X/Y, \sigma)_{fin} \longrightarrow (X/Y)_{fin}.
$$
définie ci-dessus. Alors
\begin{enumerate}
\item $t$ est représentable par des immersions ouvertes,
\item si $f$ est séparé, alors $t$ est représentable par des immersions
ouvertes et fermées,
\item si $(X/Y)_{fin}$ est un espace algébrique, alors
$(X/Y, \sigma)_{fin}$ est un espace algébrique et
un sous-espace ouvert de $(X/Y)_{fin}$, et
\item si $(X/Y)_{fin}$ est un schéma, alors $(X/Y, \sigma)_{fin}$ en est un
sous-schéma ouvert.
\end{enumerate}
\end{lemma}

\begin{proof}
Démonstration omise. Indication : étant donné un couple $(a, Z)$ au-dessus de $T$ comme dans
(\ref{equation-finite-conditions}), l'image inverse de
$Z$ par $(1_T, \sigma \circ a) : T \to T \times_Y X$ est le sous-schéma ouvert
de $T$ recherché.
\end{proof}





\section{Collections finies de flèches}
\label{section-finite-set-arrows}

\noindent
Soit $\mathcal{C}$ un groupoïde, voir
Catégories, définition \ref{categories-definition-groupoid}.
Comme expliqué dans
Groupoïdes, section \ref{groupoids-section-groupoids},
cela correspond à un septuple $(\text{Ob}, \text{Flèches}, s, t, c, e, i)$.

\medskip\noindent
À partir de ces données, nous pouvons construire un autre groupoïde
$\mathcal{C}_{fin}$ comme suit :
\begin{enumerate}
\item Un objet de $\mathcal{C}_{fin}$ consiste en une partie finie
$Z \subset \text{Flèches}$ possédant les propriétés suivantes :
\begin{enumerate}
\item $s(Z) = \{u\}$ est un singleton, et
\item $e(u) \in Z$.
\end{enumerate}
\item Un morphisme de $\mathcal{C}_{fin}$ consiste en un couple
$(Z, z)$, où $Z$ est un objet de $\mathcal{C}_{fin}$ et
$z \in Z$.
\item La source de $(Z, z)$ est $Z$.
\item Le but de $(Z, z)$ est $t(Z, z) = \{z' \circ z^{-1}; z' \in Z\}$.
\item Étant donnés $(Z_1, z_1)$ et $(Z_2, z_2)$ tels que
$s(Z_1, z_1) = t(Z_2, z_2)$, la composée $(Z_1, z_1) \circ (Z_2, z_2)$ est $(Z_2, z_1 \circ z_2)$.
\end{enumerate}
Nous omettons la vérification que ceci définit un groupoïde.
Graphiquement, un objet de $\mathcal{C}_{fin}$ peut se représenter
par un diagramme
$$
\xymatrix{
& \bullet \\
\bullet \ar@(ul, dl)[]_e \ar[ru] \ar[r] \ar[rd] & \bullet \\
& \bullet
}
$$
Pour construire un morphisme de $\mathcal{C}_{fin}$, on choisit l'une des
flèches et l'on précompose les autres flèches par son inverse. Par exemple,
si l'on choisit la flèche horizontale du milieu, le but est représenté par
$$
\xymatrix{
& \bullet \\
\bullet & \bullet \ar[l] \ar[u] \ar@(dr, ur)[]_e \ar[d] \\
& \bullet
}
$$
Notons que les cardinalités de $s(Z, z)$ et de $t(Z, z)$ sont égales.
Ainsi, $\mathcal{C}_{fin}$ est en fait une réunion disjointe dénombrable
de groupoïdes.




\section{La partie finie d'un groupoïde}
\label{section-finite-part-groupoid}

\noindent
Dans cette section, nous allons utiliser l'idée expliquée dans
la section \ref{section-finite-set-arrows}
pour prendre la partie finie d'un groupoïde en espaces algébriques.

\medskip\noindent
Soit $S$ un schéma.
Soit $B$ un espace algébrique au-dessus de $S$.
Soit $(U, R, s, t, c, e, i)$ un groupoïde en espaces algébriques au-dessus de $B$.
Hypothèse : les morphismes $s, t$ sont séparés et localement de type fini.
Cette notation et cette hypothèse resteront fixées dans toute la section.

\medskip\noindent
Notons $R_s$ l'espace algébrique $R$ considéré comme un
espace algébrique au-dessus de $U$ par l'intermédiaire de $s$. Posons
$U' = (R_s/U, e)_{fin}$. Puisque $s$ est séparé et localement
de type fini, la
proposition \ref{proposition-finite-algebraic-space} et le
lemme \ref{lemma-finite-plus-section}
montrent que $U'$ est un espace algébrique muni d'un morphisme \'etale
$g : U' \to U$. De plus, d'après le
lemme \ref{lemma-finite-sheaf},
il existe un sous-espace ouvert universel
$Z_{univ} \subset R \times_{s, U, g} U'$, fini au-dessus de $U'$,
tel que $(1_{U'}, e \circ g) : U' \to R \times_{s, U, g} U'$
se factorise par $Z_{univ}$. En outre, d'après le
lemme \ref{lemma-finite-separated},
le sous-espace ouvert $Z_{univ}$ est aussi fermé dans
$R \times_{s, U, g} U'$. Voici la situation obtenue :
$$
\xymatrix{
Z_{univ} \ar[d] \ar[rd] & \\
R \times_{s, U, g} U' \ar[d] \ar[r] & U' \ar[d]^g \\
R \ar[r]^s & U
}
$$
Soit $T$ un schéma au-dessus de $B$. On voit qu'un point de
$Z_{univ}$ à valeurs dans $T$ peut être considéré comme un triplet
$(u, Z, z)$ où
\begin{enumerate}
\item $u : T \to U$ est un point de $U$ à valeurs dans $T$,
\item $Z \subset R \times_{s, U, u} T$ est un sous-espace ouvert et fermé,
fini au-dessus de $T$, par lequel $(e \circ u, 1_T)$ se factorise, et
\item $z : T \to R$ est un point de $R$ à valeurs dans $T$ tel que
$s \circ z = u$ et tel que $(z, 1_T)$ se factorise par $Z$.
\end{enumerate}
Cela étant dit, la discussion de la
section \ref{section-finite-set-arrows}
montre moralement que l'on peut munir $(Z_{univ}, U')$ d'une structure de
groupoïde en espaces algébriques au-dessus de $B$. Pour nous en assurer,
définissons un à un les morphismes $s', t', c', e', i'$ du point de vue
fonctoriel. (Nous invitons le lecteur à ne pas lire ce qui suit avant d'avoir
lu et compris la construction simple de la
section \ref{section-finite-set-arrows}.)

\medskip\noindent
Le morphisme $s' : Z_{univ} \to U'$ correspond à la règle
$$
s' : (u, Z, z) \mapsto (u, Z).
$$
Le morphisme $t' : Z_{univ} \to U'$ est donné par la règle
$$
t' : (u, Z, z) \mapsto (t \circ z, c(Z, i \circ z)).
$$
Le terme $c(Z, i \circ z)$
a un sens, car l'application
$c(-, i \circ z) : R \times_{s, U, u} T \to R \times_{s, U, t \circ z} T$
est un isomorphisme d'inverse $c(-, z)$.
Le morphisme $e' : U' \to Z_{univ}$ est donné par la règle
$$
e' : (u, Z) \mapsto (u, Z, (e \circ u, 1_T)).
$$
Notons que ceci a un sens en vertu de la condition imposant que
$(e \circ u, 1_T)$ se factorise par $Z$.
Le morphisme $i' : Z_{univ} \to Z_{univ}$ est donné par la règle
$$
i' : (u, Z, z) \mapsto (t \circ z, c(Z, i \circ z), i \circ z).
$$
Enfin, la composition est définie par la règle
$$
c' : ((u_1, Z_1, z_1), (u_2, Z_2, z_2)) \mapsto (u_2, Z_2, z_1 \circ z_2).
$$
Nous omettons la vérification que les axiomes d'un groupoïde en espaces
algébriques sont satisfaits par $(U', Z_{univ}, s', t', c', e', i')$.

\medskip\noindent
Enfin, il existe un morphisme canonique de groupoïdes
$$
(U', Z_{univ}, s', t', c', e', i')
\longrightarrow
(U, R, s, t, c, e, i)
$$
En effet, le morphisme $U' \to U$ est le morphisme $g : U' \to U$,
qui est défini par la règle $(u, Z) \mapsto u$. Le morphisme
$Z_{univ} \to R$ est défini par la règle $(u, Z, z) \mapsto z$.
Ceci achève la construction. Résumons les résultats obtenus.

\begin{lemma}
\label{lemma-finite-part-groupoid}
Soit $S$ un schéma.
Soit $B$ un espace algébrique au-dessus de $S$.
Soit $(U, R, s, t, c, e, i)$ un groupoïde en espaces algébriques au-dessus de $B$.
Supposons que les morphismes $s, t$ soient séparés et localement de type fini.
Il existe un morphisme canonique
$$
(U', Z_{univ}, s', t', c', e', i')
\longrightarrow
(U, R, s, t, c, e, i)
$$
de groupoïdes en espaces algébriques au-dessus de $B$, où
\begin{enumerate}
\item $g : U' \to U$ s'identifie à $(R_s/U, e)_{fin} \to U$, et
\item $Z_{univ} \subset R \times_{s, U, g} U'$ est le sous-espace
ouvert (et fermé) universel, fini au-dessus de $U'$, qui contient le
changement de base de l'unité $e$.
\end{enumerate}
\end{lemma}

\begin{proof}
Voir la discussion ci-dessus.
\end{proof}









\section{Localisation \'etale des schémas en groupoïdes}
\label{section-etale-localize}

\noindent
Dans cette section, nous démontrons des résultats analogues à
\cite[Proposition 4.2]{K-M}. Nous cherchons à être un peu plus généraux et
à éviter l'emploi des schémas de Hilbert en utilisant plutôt la partie finie
d'un morphisme. Le but est de « scinder » un groupoïde en espaces algébriques
au-dessus d'un point après localisation \'etale. Voici la définition
(très proche de \cite[définition 4.1]{K-M}).

\begin{definition}
\label{definition-split-at-point}
Soit $S$ un schéma. Soit $B$ un espace algébrique au-dessus de $S$.
Soit $(U, R, s, t, c)$ un groupoïde en espaces algébriques au-dessus de $B$.
Soit $u \in |U|$ un point.
\begin{enumerate}
\item On dit que $R$ est {\it fortement scindé au-dessus de $u$} s'il existe
un sous-espace ouvert $P \subset R$ tel que
\begin{enumerate}
\item $(U, P, s|_P, t|_P, c|_{P \times_{s, U, t} P})$ est un
groupoïde en espaces algébriques au-dessus de $B$,
\item $s|_P$, $t|_P$ sont finis, et
\item $\{r \in |R| : s(r) = u, t(r) = u\} \subset |P|$.
\end{enumerate}
Le choix d'un tel $P$ sera appelé un
{\it scindage fort de $R$ au-dessus de $u$}.
\item On dit que $R$ est {\it scindé au-dessus de $u$} s'il existe un
sous-espace ouvert $P \subset R$ tel que
\begin{enumerate}
\item $(U, P, s|_P, t|_P, c|_{P \times_{s, U, t} P})$ est un
groupoïde en espaces algébriques au-dessus de $B$,
\item $s|_P$, $t|_P$ sont finis, et
\item $\{g \in |G| : g\text{ est au-dessus de }u\} \subset |P|$, où
$G \to U$ est le stabilisateur.
\end{enumerate}
Le choix d'un tel $P$ sera appelé un
{\it scindage de $R$ au-dessus de $u$}.
\item On dit que $R$ est {\it quasi-scindé au-dessus de $u$} s'il existe
un sous-espace ouvert $P \subset R$ tel que
\begin{enumerate}
\item $(U, P, s|_P, t|_P, c|_{P \times_{s, U, t} P})$ est un
groupoïde en espaces algébriques au-dessus de $B$,
\item $s|_P$, $t|_P$ sont finis, et
\item $e(u) \in |P|$\footnote{Cette condition découle de (a).}.
\end{enumerate}
Le choix d'un tel $P$ sera appelé un {\it quasi-scindage de $R$ au-dessus de $u$}.
\end{enumerate}
\end{definition}

\noindent
Notons la similitude entre les conditions imposées à $P$ et celles imposées
aux couples dans (\ref{equation-finite-conditions}). En particulier, si
$s, t$ sont séparés, alors $P$ est aussi fermé dans $R$ (voir le
lemme \ref{lemma-finite-separated}).

\medskip\noindent
Supposons donné un groupoïde en espaces algébriques
$(U, R, s, t, c)$ au-dessus de $B$ et un point $u \in |U|$.
Puisque le but est de scinder le groupoïde après localisation \'etale,
on peut tout aussi bien remplacer $U$ par un schéma affine (nous voulons dire
que cela est sans conséquence pour toute application éventuelle).
En outre, les hypothèses supplémentaires que nous devrons
imposer feront de $R$ un schéma, au moins dans un voisinage
de $\{r \in |R| : s(r) = u, t(r) = u\}$ ou de $e(u)$. C'est pourquoi
nous partons d'un schéma en groupoïdes comme décrit ci-dessous.
Cependant, notre méthode de démonstration nous fait sortir de la catégorie
des schémas, raison pour laquelle nous avons formulé ci-dessus le scindage
pour les groupoïdes en espaces algébriques.
D'autre part, nous ne connaissons d'applications que
dans le cas où les morphismes $s$, $t$
sont aussi plats et de présentation finie ; dans ce cas,
nous revenons finalement dans la catégorie des schémas.

\begin{situation}[Scindage fort]
\label{situation-strong-splitting}
Soit $S$ un schéma.
Soit $(U, R, s, t, c)$ un schéma en groupoïdes au-dessus de $S$.
Soit $u \in U$ un point. Supposons que
\begin{enumerate}
\item $s, t : R \to U$ soient séparés,
\item $s$, $t$ soient localement de type fini,
\item l'ensemble $\{r \in R : s(r) = u, t(r) = u\}$ soit fini, et
\item $s$ soit quasi-fini en chaque point de l'ensemble de (3).
\end{enumerate}
Notons que les hypothèses (3) et (4) découlent de l'hypothèse
que la fibre $s^{-1}(\{u\})$ est finie, voir
Morphismes, lemme \ref{morphisms-lemma-finite-fibre}.
\end{situation}

\begin{situation}[Scindage]
\label{situation-splitting}
Soit $S$ un schéma.
Soit $(U, R, s, t, c)$ un schéma en groupoïdes au-dessus de $S$.
Soit $u \in U$ un point. Supposons que
\begin{enumerate}
\item $s, t : R \to U$ soient séparés,
\item $s$, $t$ soient localement de type fini,
\item l'ensemble $\{g \in G : g\text{ est au-dessus de }u\}$ soit fini, où $G \to U$
est le stabilisateur, et
\item $s$ soit quasi-fini en chaque point de l'ensemble de (3).
\end{enumerate}
\end{situation}

\begin{situation}[Quasi-scindage]
\label{situation-quasi-splitting}
Soit $S$ un schéma.
Soit $(U, R, s, t, c)$ un schéma en groupoïdes au-dessus de $S$.
Soit $u \in U$ un point. Supposons que
\begin{enumerate}
\item $s, t : R \to U$ soient séparés,
\item $s$, $t$ soient localement de type fini, et
\item $s$ soit quasi-fini en $e(u)$.
\end{enumerate}
\end{situation}

\noindent
Pour notre application aux théorèmes d'existence des espaces algébriques,
le cas des quasi-scindages suffit. En outre, le cas du quasi-scindage
nous permettra de démontrer, pour les champs quasi-DM, un théorème de
structure locale pour la topologie \'etale. Le cas du scindage servira à démontrer une version
du théorème de Keel--Mori. Le cas du scindage fort fournit
un théorème de structure locale pour la topologie \'etale des champs
algébriques quasi-DM à diagonale quasi-compacte.

\begin{lemma}[Existence d'un scindage fort]
\label{lemma-strong-splitting}
Dans la situation \ref{situation-strong-splitting},
il existe un espace algébrique $U'$, un morphisme \'etale
$U' \to U$ et un point $u' : \Spec(\kappa(u)) \to U'$
au-dessus de $u : \Spec(\kappa(u)) \to U$,
tels que la restriction $R' = R|_{U'}$ de $R$ à $U'$
soit fortement scindée au-dessus de $u'$.
\end{lemma}

\begin{proof}
Soit $f : (U', Z_{univ}, s', t', c') \to (U, R, s, t, c)$ construit dans le
lemme \ref{lemma-finite-part-groupoid}.
Rappelons que $R' = R \times_{(U \times_S U)} (U' \times_S U')$.
On obtient ainsi un morphisme $(f, t', s') : Z_{univ} \to R'$ de groupoïdes
en espaces algébriques
$$
(U', Z_{univ}, s', t', c') \to (U', R', s', t', c')
$$
(par abus de notation, nous désignons les morphismes des deux groupoïdes
par les mêmes symboles). Or $Z_{univ} \subset R \times_{s, U, g} U'$ est ouvert,
et $R' \to R \times_{s, U, g} U'$ est \'etale (comme changement de base
de $U' \to U$) ; ainsi, $Z_{univ} \to R'$ est une immersion ouverte.
Par construction, les morphismes $s', t' : Z_{univ} \to U'$ sont finis.
Il reste à trouver le point $u'$ de $U'$.

\medskip\noindent
Considérons $u$ comme un morphisme $\Spec(\kappa(u)) \to U$, comme dans
l'énoncé du lemme. Posons $F_u = R \times_{s, U} \Spec(\kappa(u))$.
L'ensemble $\{r \in R : s(r) = u, t(r) = u\}$ est fini par hypothèse
et $F_u \to \Spec(\kappa(u))$ est quasi-fini en chacun
de ses éléments par hypothèse. On peut donc trouver une décomposition en
sous-schémas ouverts et fermés
$$
F_u = Z_u \amalg Reste
$$
où $Z_u$ est un schéma fini au-dessus de $\kappa(u)$ dont le support est
$\{r \in R : s(r) = u, t(r) = u\}$. Notons que $e(u) \in Z_u$.
Par conséquent, d'après la construction de $U'$ dans la
section \ref{section-finite-part-groupoid},
$(u, Z_u)$ définit un point de $U'$ à valeurs dans
$\Spec(\kappa(u))$, que nous notons $u'$.

\medskip\noindent
Il nous reste à montrer que l'ensemble
$\{r' \in |R'| : s'(r') = u', t'(r') = u'\}$
est contenu dans $|Z_{univ}|$.
Choisissons un point quelconque $r'$ de cet ensemble et représentons-le par
un morphisme $z' : \Spec(k) \to R'$. Notons $z : \Spec(k) \to R$
la composée de $z'$ et de l'application $R' \to R$.
Il est clair que $z$ définit un élément de l'ensemble
$\{r \in R : s(r) = u, t(r) = u\}$.
De plus, les composées $s \circ z, t \circ z : \Spec(k) \to U$
se factorisent par $u$ ; nous pouvons donc considérer $s \circ z, t \circ z$
comme des morphismes $\Spec(k) \to \Spec(\kappa(u))$. Alors
$z' = (z, u' \circ t \circ z, u'\circ s \circ z)$ comme morphismes vers
$R' = R \times_{(U \times_S U)} (U' \times_S U')$.
Considérons le triplet
$$
(s \circ z, Z_u \times_{\Spec(\kappa(u)), s \circ z} \Spec(k), z)
$$
où $Z_u$ est défini comme ci-dessus. Ceci définit un point de $Z_{univ}$
à valeurs dans $\Spec(k)$, dont l'image dans $U'$ par $s', t'$ est $u'$ et
dont l'image par
$Z_{univ} \to R'$ est le point $r'$, en vertu de la relation entre
$z$ et $z'$ mentionnée ci-dessus.
Ceci achève la démonstration.
\end{proof}

\begin{lemma}[Existence d'un scindage]
\label{lemma-splitting}
Dans la situation \ref{situation-splitting},
il existe un espace algébrique $U'$, un morphisme \'etale
$U' \to U$ et un point $u' : \Spec(\kappa(u)) \to U'$
au-dessus de $u : \Spec(\kappa(u)) \to U$,
tels que la restriction $R' = R|_{U'}$ de $R$ à $U'$
soit scindée au-dessus de $u'$.
\end{lemma}

\begin{proof}
Soit $f : (U', Z_{univ}, s', t', c') \to (U, R, s, t, c)$ construit dans le
lemme \ref{lemma-finite-part-groupoid}.
Rappelons que $R' = R \times_{(U \times_S U)} (U' \times_S U')$.
On obtient ainsi un morphisme $(f, t', s') : Z_{univ} \to R'$ de groupoïdes
en espaces algébriques
$$
(U', Z_{univ}, s', t', c') \to (U', R', s', t', c')
$$
(par abus de notation, nous désignons les morphismes des deux groupoïdes
par les mêmes symboles). Or $Z_{univ} \subset R \times_{s, U, g} U'$ est ouvert,
et $R' \to R \times_{s, U, g} U'$ est \'etale (comme changement de base
de $U' \to U$) ; ainsi, $Z_{univ} \to R'$ est une immersion ouverte.
Par construction, les morphismes $s', t' : Z_{univ} \to U'$ sont finis.
Il reste à trouver le point $u'$ de $U'$.

\medskip\noindent
Considérons $u$ comme un morphisme $\Spec(\kappa(u)) \to U$, comme dans
l'énoncé du lemme. Posons $F_u = R \times_{s, U} \Spec(\kappa(u))$.
Soit $G_u \subset F_u$ la fibre schématique de $G \to U$ au-dessus de $u$.
Par hypothèse, $G_u$ est fini et $F_u \to \Spec(\kappa(u))$
est quasi-fini en chaque point de $G_u$.
On peut donc trouver une décomposition en sous-schémas ouverts et fermés
$$
F_u = Z_u \amalg Reste
$$
où $Z_u$ est un schéma fini au-dessus de $\kappa(u)$ dont le support est $G_u$.
Notons que $e(u) \in Z_u$. Par conséquent, d'après la construction de $U'$ dans la
section \ref{section-finite-part-groupoid},
$(u, Z_u)$ définit un point de $U'$ à valeurs dans
$\Spec(\kappa(u))$, que nous notons $u'$.

\medskip\noindent
Il nous reste à montrer que l'ensemble $\{g' \in |G'| : g'\text{ est au-dessus de }u'\}$
est contenu dans $|Z_{univ}|$. Choisissons un point $g'$ de cet ensemble et
représentons-le par un morphisme $z' : \Spec(k) \to G'$. Notons
$z : \Spec(k) \to G$ la composée de $z'$ et de l'application $G' \to G$.
Il est clair que $z$ définit un point de $G_u$. Écrivons plus précisément
$\tilde u : \Spec(k) \to u \to U$ pour l'application correspondante vers $u$ ou $U$.
Considérons le triplet
$$
(\tilde u, Z_u \times_{u, \tilde u} \Spec(k), z)
$$
où $Z_u$ est défini comme ci-dessus. Ceci définit un point de $Z_{univ}$
à valeurs dans $\Spec(k)$, dont l'image dans $U'$ par $s', t'$ est $u'$ et
dont l'image par $Z_{univ} \to R'$ est le point $z'$
(car son image dans $R$ est $z$).
Ceci achève la démonstration.
\end{proof}

\begin{lemma}[Existence d'un quasi-scindage]
\label{lemma-quasi-splitting}
Dans la situation \ref{situation-quasi-splitting},
il existe un espace algébrique $U'$, un morphisme \'etale
$U' \to U$ et un point $u' : \Spec(\kappa(u)) \to U'$
au-dessus de $u : \Spec(\kappa(u)) \to U$,
tels que la restriction $R' = R|_{U'}$ de $R$ à $U'$
soit quasi-scindée au-dessus de $u'$.
\end{lemma}

\begin{proof}
Soit $f : (U', Z_{univ}, s', t', c') \to (U, R, s, t, c)$ construit dans le
lemme \ref{lemma-finite-part-groupoid}.
Rappelons que $R' = R \times_{(U \times_S U)} (U' \times_S U')$.
On obtient ainsi un morphisme $(f, t', s') : Z_{univ} \to R'$ de groupoïdes
en espaces algébriques
$$
(U', Z_{univ}, s', t', c') \to (U', R', s', t', c')
$$
(par abus de notation, nous désignons les morphismes des deux groupoïdes
par les mêmes symboles). Or $Z_{univ} \subset R \times_{s, U, g} U'$ est ouvert,
et $R' \to R \times_{s, U, g} U'$ est \'etale (comme changement de base
de $U' \to U$) ; ainsi, $Z_{univ} \to R'$ est une immersion ouverte.
Par construction, les morphismes $s', t' : Z_{univ} \to U'$ sont finis.
Il reste à trouver le point $u'$ de $U'$.

\medskip\noindent
Considérons $u$ comme un morphisme $\Spec(\kappa(u)) \to U$, comme dans
l'énoncé du lemme. Posons $F_u = R \times_{s, U} \Spec(\kappa(u))$.
Le morphisme $F_u \to \Spec(\kappa(u))$ est quasi-fini en $e(u)$
par hypothèse. On peut donc trouver une décomposition en sous-schémas
ouverts et fermés
$$
F_u = Z_u \amalg Reste
$$
où $Z_u$ est un schéma fini au-dessus de $\kappa(u)$ dont le support est $e(u)$.
Par conséquent, d'après la construction de $U'$ dans la
section \ref{section-finite-part-groupoid},
$(u, Z_u)$ définit un point de $U'$ à valeurs dans
$\Spec(\kappa(u))$, que nous notons $u'$. Pour achever la démonstration,
il suffit de montrer que $e'(u') \in Z_{univ}$, ce qui est clair.
\end{proof}

\noindent
Enfin, sous des hypothèses supplémentaires, nous obtenons des schémas.

\begin{lemma}
\label{lemma-strong-splitting-scheme}
Dans la situation \ref{situation-strong-splitting}, supposons en outre que
$s, t$ soient plats et localement de présentation finie.
Alors il existe un schéma $U'$, un morphisme \'etale séparé
$U' \to U$ et un point $u' \in U'$
au-dessus de $u$, avec $\kappa(u) = \kappa(u')$,
tels que la restriction $R' = R|_{U'}$ de $R$ à $U'$
soit fortement scindée au-dessus de $u'$.
\end{lemma}

\begin{proof}
Ceci résulte de la construction de $U'$ dans la démonstration du
lemme \ref{lemma-strong-splitting},
car, dans ce cas, $U' = (R_s/U, e)_{fin}$ est un schéma séparé au-dessus
de $U$ d'après les
lemmes \ref{lemma-finite-separated-flat-locally-finite-presentation} et
\ref{lemma-finite-plus-section}.
\end{proof}

\begin{lemma}
\label{lemma-splitting-scheme}
Dans la situation \ref{situation-splitting}, supposons en outre que
$s, t$ soient plats et localement de présentation finie.
Alors il existe un schéma $U'$, un morphisme \'etale séparé
$U' \to U$ et un point $u' \in U'$
au-dessus de $u$, avec $\kappa(u) = \kappa(u')$,
tels que la restriction $R' = R|_{U'}$ de $R$ à $U'$
soit scindée au-dessus de $u'$.
\end{lemma}

\begin{proof}
Ceci résulte de la construction de $U'$ dans la démonstration du
lemme \ref{lemma-splitting},
car, dans ce cas, $U' = (R_s/U, e)_{fin}$ est un schéma séparé au-dessus
de $U$ d'après les
lemmes \ref{lemma-finite-separated-flat-locally-finite-presentation} et
\ref{lemma-finite-plus-section}.
\end{proof}

\begin{lemma}
\label{lemma-quasi-splitting-scheme}
Dans la situation \ref{situation-quasi-splitting}, supposons en outre que
$s, t$ soient plats et localement de présentation finie.
Alors il existe un schéma $U'$, un morphisme \'etale séparé
$U' \to U$ et un point $u' \in U'$ au-dessus de $u$, avec
$\kappa(u) = \kappa(u')$, tels que la restriction $R' = R|_{U'}$ de
$R$ à $U'$ soit quasi-scindée au-dessus de $u'$.
\end{lemma}

\begin{proof}
Ceci résulte de la construction de $U'$ dans la démonstration du
lemme \ref{lemma-quasi-splitting},
car, dans ce cas, $U' = (R_s/U, e)_{fin}$ est un schéma séparé
au-dessus de $U$ d'après les
lemmes \ref{lemma-finite-separated-flat-locally-finite-presentation} et
\ref{lemma-finite-plus-section}.
\end{proof}

\noindent
En fait, nous pouvons obtenir des schémas affines en appliquant un résultat
antérieur sur les groupoïdes finis localement libres.

\begin{lemma}
\label{lemma-strong-splitting-affine-scheme}
Dans la situation \ref{situation-strong-splitting}, supposons en outre que
$s, t$ soient plats et localement de présentation finie et que $U$ soit affine.
Alors il existe un schéma affine $U'$, un morphisme \'etale
$U' \to U$ et un point $u' \in U'$ au-dessus de $u$, avec
$\kappa(u) = \kappa(u')$, tels que la restriction $R' = R|_{U'}$ de
$R$ à $U'$ soit fortement scindée au-dessus de $u'$.
\end{lemma}

\begin{proof}
Soient $U' \to U$ le morphisme \'etale séparé de schémas et $u' \in U'$ le point
obtenus dans le lemme \ref{lemma-strong-splitting-scheme}.
Soit $P \subset R'$ le scindage fort de $R'$ au-dessus de $u'$. D'après
Compléments sur les groupoïdes, lemme \ref{more-groupoids-lemma-restrict-preserves-type},
les morphismes $s'|_P, t'|_P : P \to U'$ sont plats et localement de présentation finie.
Ils sont finis par définition du scindage. Ainsi, $s'|_P, t'|_P$ sont finis
localement libres, voir
Morphismes, lemme \ref{morphisms-lemma-finite-flat}.
En particulier, $(t'|_P)((s'|_P)^{-1}(u'))$ est un ensemble fini de points
$\{u'_1, u'_2, \ldots, u'_n\}$ de $U'$. Choisissons un ouvert quasi-compact
$W \subset U'$ contenant chaque $u'_i$. Puisque $U$ est affine, le morphisme
$W \to U$ est quasi-compact (voir
Schémas, lemme \ref{schemes-lemma-quasi-compact-affine}).
Le morphisme $W \to U$ est aussi localement quasi-fini (voir
Morphismes, lemme \ref{morphisms-lemma-etale-locally-quasi-finite})
et séparé. Par conséquent, d'après
Compléments sur les morphismes,
lemme \ref{more-morphisms-lemma-quasi-finite-separated-quasi-affine}
(une version du théorème principal de Zariski),
nous concluons que $W$ est quasi-affine. D'après
Propriétés, lemme \ref{properties-lemma-ample-finite-set-in-affine},
l'ensemble $\{u'_1, \ldots, u'_n\}$ est contenu dans un ouvert affine
de $U'$. Nous pouvons donc appliquer
Groupoïdes, lemme \ref{groupoids-lemma-find-invariant-affine}
pour conclure qu'il existe un ouvert affine $U'' \subset U'$, invariant par
$P$, qui contient $u'$.

\medskip\noindent
Pour achever la démonstration, notons $R'' = R|_{U''}$ la restriction de $R$
à $U''$. C'est aussi la restriction de $R'$ à $U''$.
Puisque $P \subset R'$ est un sous-schéma ouvert et fermé, il en va de même de
$P|_{U''} \subset R''$. Par construction, le sous-schéma ouvert $U'' \subset U'$
est invariant par $P$, ce qui signifie que
$P|_{U''} = (s'|_P)^{-1}(U'') = (t'|_P)^{-1}(U'')$
(voir la discussion dans
Groupoïdes, section \ref{groupoids-section-invariant}) ;
les restrictions de $s''$ et de $t''$ à $P|_{U''}$ sont donc encore finies.
Le sous-schéma en groupoïdes $P|_{U''}$ est encore un scindage fort de
$R''$ au-dessus de $u'$ : les conditions (a) et (b) ont été vérifiées ci-dessus,
et (c) résulte immédiatement de $\{r' \in R': t'(r') = u', s'(r') = u'\} =
\{r'' \in R'': t''(r'') = u', s''(r'') = u'\}$.
Le lemme est démontré.
\end{proof}

\begin{lemma}
\label{lemma-splitting-affine-scheme}
Dans la situation \ref{situation-splitting}, supposons en outre que
$s, t$ soient plats et localement de présentation finie et que $U$ soit affine.
Alors il existe un schéma affine $U'$, un morphisme \'etale
$U' \to U$ et un point $u' \in U'$ au-dessus de $u$, avec
$\kappa(u) = \kappa(u')$, tels que la restriction $R' = R|_{U'}$ de
$R$ à $U'$ soit scindée au-dessus de $u'$.
\end{lemma}

\begin{proof}
La démonstration de ce lemme est littéralement la même que celle du
lemme \ref{lemma-strong-splitting-affine-scheme},
si ce n'est que « scindage fort » doit être remplacé par « scindage »
(deux fois) et que le renvoi au
lemme \ref{lemma-strong-splitting-scheme}
doit être remplacé par un renvoi au
lemme \ref{lemma-splitting-scheme}.
\end{proof}

\begin{lemma}
\label{lemma-quasi-splitting-affine-scheme}
Dans la situation \ref{situation-quasi-splitting}, supposons en outre que
$s, t$ soient plats et localement de présentation finie et que $U$ soit affine.
Alors il existe un schéma affine $U'$, un morphisme \'etale
$U' \to U$ et un point $u' \in U'$ au-dessus de $u$, avec
$\kappa(u) = \kappa(u')$, tels que la restriction $R' = R|_{U'}$ de
$R$ à $U'$ soit quasi-scindée au-dessus de $u'$.
\end{lemma}

\begin{proof}
La démonstration de ce lemme est littéralement la même que celle du
lemme \ref{lemma-strong-splitting-affine-scheme},
si ce n'est que « scindage fort » doit être remplacé par « quasi-scindage »
(deux fois) et que le renvoi au
lemme \ref{lemma-strong-splitting-scheme}
doit être remplacé par un renvoi au
lemme \ref{lemma-quasi-splitting-scheme}.
\end{proof}


\begin{multicols}{2}[\section{Autres chapitres}]
\noindent
Préliminaires
\begin{enumerate}
\item \hyperref[introduction-section-phantom]{Introduction}
\item \hyperref[conventions-section-phantom]{Conventions}
\item \hyperref[sets-section-phantom]{Théorie des ensembles}
\item \hyperref[categories-section-phantom]{Catégories}
\item \hyperref[topology-section-phantom]{Topologie}
\item \hyperref[sheaves-section-phantom]{Faisceaux sur les espaces}
\item \hyperref[sites-section-phantom]{Sites et faisceaux}
\item \hyperref[stacks-section-phantom]{Champs}
\item \hyperref[fields-section-phantom]{Corps}
\item \hyperref[algebra-section-phantom]{Algèbre commutative}
\item \hyperref[brauer-section-phantom]{Groupes de Brauer}
\item \hyperref[homology-section-phantom]{Algèbre homologique}
\item \hyperref[derived-section-phantom]{Catégories dérivées}
\item \hyperref[simplicial-section-phantom]{Méthodes simpliciales}
\item \hyperref[more-algebra-section-phantom]{Compléments d'algèbre}
\item \hyperref[smoothing-section-phantom]{Lissage des morphismes d'anneaux}
\item \hyperref[modules-section-phantom]{Faisceaux de modules}
\item \hyperref[sites-modules-section-phantom]{Modules sur les sites}
\item \hyperref[injectives-section-phantom]{Objets injectifs}
\item \hyperref[cohomology-section-phantom]{Cohomologie des faisceaux}
\item \hyperref[sites-cohomology-section-phantom]{Cohomologie sur les sites}
\item \hyperref[dga-section-phantom]{Algèbre différentielle graduée}
\item \hyperref[dpa-section-phantom]{Algèbre à puissances divisées}
\item \hyperref[sdga-section-phantom]{Faisceaux différentiels gradués}
\item \hyperref[hypercovering-section-phantom]{Hyperrecouvrements}
\end{enumerate}
Schémas
\begin{enumerate}
\setcounter{enumi}{25}
\item \hyperref[schemes-section-phantom]{Schémas}
\item \hyperref[constructions-section-phantom]{Constructions de schémas}
\item \hyperref[properties-section-phantom]{Propriétés des schémas}
\item \hyperref[morphisms-section-phantom]{Morphismes de schémas}
\item \hyperref[coherent-section-phantom]{Cohomologie des schémas}
\item \hyperref[divisors-section-phantom]{Diviseurs}
\item \hyperref[limits-section-phantom]{Limites de schémas}
\item \hyperref[varieties-section-phantom]{Variétés}
\item \hyperref[topologies-section-phantom]{Topologies sur les schémas}
\item \hyperref[descent-section-phantom]{Descente}
\item \hyperref[perfect-section-phantom]{Catégories dérivées de schémas}
\item \hyperref[more-morphisms-section-phantom]{Compléments sur les morphismes}
\item \hyperref[flat-section-phantom]{Compléments sur la platitude}
\item \hyperref[groupoids-section-phantom]{Schémas en groupoïdes}
\item \hyperref[more-groupoids-section-phantom]{Compléments sur les schémas en groupoïdes}
\item \hyperref[etale-section-phantom]{Morphismes étales de schémas}
\end{enumerate}
Thèmes de théorie des schémas
\begin{enumerate}
\setcounter{enumi}{41}
\item \hyperref[chow-section-phantom]{Homologie de Chow}
\item \hyperref[intersection-section-phantom]{Théorie de l'intersection}
\item \hyperref[pic-section-phantom]{Schémas de Picard des courbes}
\item \hyperref[weil-section-phantom]{Théories cohomologiques de Weil}
\item \hyperref[adequate-section-phantom]{Modules adéquats}
\item \hyperref[dualizing-section-phantom]{Complexes dualisants}
\item \hyperref[duality-section-phantom]{Dualité pour les schémas}
\item \hyperref[discriminant-section-phantom]{Discriminants et différentes}
\item \hyperref[derham-section-phantom]{Cohomologie de de Rham}
\item \hyperref[local-cohomology-section-phantom]{Cohomologie locale}
\item \hyperref[algebraization-section-phantom]{Géométrie algébrique et formelle}
\item \hyperref[curves-section-phantom]{Courbes algébriques}
\item \hyperref[resolve-section-phantom]{Résolution des surfaces}
\item \hyperref[models-section-phantom]{Réduction semi-stable}
\item \hyperref[functors-section-phantom]{Foncteurs et morphismes}
\item \hyperref[equiv-section-phantom]{Catégories dérivées de variétés}
\item \hyperref[pione-section-phantom]{Groupes fondamentaux de schémas}
\item \hyperref[etale-cohomology-section-phantom]{Cohomologie étale}
\item \hyperref[crystalline-section-phantom]{Cohomologie cristalline}
\item \hyperref[proetale-section-phantom]{Cohomologie pro-étale}
\item \hyperref[relative-cycles-section-phantom]{Cycles relatifs}
\item \hyperref[more-etale-section-phantom]{Compléments de cohomologie étale}
\item \hyperref[trace-section-phantom]{La formule des traces}
\end{enumerate}
Espaces algébriques
\begin{enumerate}
\setcounter{enumi}{64}
\item \hyperref[spaces-section-phantom]{Espaces algébriques}
\item \hyperref[spaces-properties-section-phantom]{Propriétés des espaces algébriques}
\item \hyperref[spaces-morphisms-section-phantom]{Morphismes d'espaces algébriques}
\item \hyperref[decent-spaces-section-phantom]{Espaces algébriques décents}
\item \hyperref[spaces-cohomology-section-phantom]{Cohomologie des espaces algébriques}
\item \hyperref[spaces-limits-section-phantom]{Limites d'espaces algébriques}
\item \hyperref[spaces-divisors-section-phantom]{Diviseurs sur les espaces algébriques}
\item \hyperref[spaces-over-fields-section-phantom]{Espaces algébriques sur des corps}
\item \hyperref[spaces-topologies-section-phantom]{Topologies sur les espaces algébriques}
\item \hyperref[spaces-descent-section-phantom]{Descente et espaces algébriques}
\item \hyperref[spaces-perfect-section-phantom]{Catégories dérivées d'espaces}
\item \hyperref[spaces-more-morphisms-section-phantom]{Compléments sur les morphismes d'espaces}
\item \hyperref[spaces-flat-section-phantom]{Platitude sur les espaces algébriques}
\item \hyperref[spaces-groupoids-section-phantom]{Groupoïdes dans les espaces algébriques}
\item \hyperref[spaces-more-groupoids-section-phantom]{Compléments sur les groupoïdes dans les espaces}
\item \hyperref[bootstrap-section-phantom]{Amorçage}
\item \hyperref[spaces-pushouts-section-phantom]{Sommes amalgamées d'espaces algébriques}
\end{enumerate}
Thèmes de géométrie
\begin{enumerate}
\setcounter{enumi}{81}
\item \hyperref[spaces-chow-section-phantom]{Groupes de Chow des espaces}
\item \hyperref[groupoids-quotients-section-phantom]{Quotients de groupoïdes}
\item \hyperref[spaces-more-cohomology-section-phantom]{Compléments sur la cohomologie des espaces}
\item \hyperref[spaces-simplicial-section-phantom]{Espaces simpliciaux}
\item \hyperref[spaces-duality-section-phantom]{Dualité pour les espaces}
\item \hyperref[formal-spaces-section-phantom]{Espaces algébriques formels}
\item \hyperref[restricted-section-phantom]{Algébrisation des espaces formels}
\item \hyperref[spaces-resolve-section-phantom]{Résolution des surfaces revisitée}
\end{enumerate}
Théorie des déformations
\begin{enumerate}
\setcounter{enumi}{89}
\item \hyperref[formal-defos-section-phantom]{Théorie formelle des déformations}
\item \hyperref[defos-section-phantom]{Théorie des déformations}
\item \hyperref[cotangent-section-phantom]{Le complexe cotangent}
\item \hyperref[examples-defos-section-phantom]{Problèmes de déformation}
\end{enumerate}
Champs algébriques
\begin{enumerate}
\setcounter{enumi}{93}
\item \hyperref[algebraic-section-phantom]{Champs algébriques}
\item \hyperref[examples-stacks-section-phantom]{Exemples de champs}
\item \hyperref[stacks-sheaves-section-phantom]{Faisceaux sur les champs algébriques}
\item \hyperref[criteria-section-phantom]{Critères de représentabilité}
\item \hyperref[artin-section-phantom]{Axiomes d'Artin}
\item \hyperref[quot-section-phantom]{Espaces de Quot et de Hilbert}
\item \hyperref[stacks-properties-section-phantom]{Propriétés des champs algébriques}
\item \hyperref[stacks-morphisms-section-phantom]{Morphismes de champs algébriques}
\item \hyperref[stacks-limits-section-phantom]{Limites de champs algébriques}
\item \hyperref[stacks-cohomology-section-phantom]{Cohomologie des champs algébriques}
\item \hyperref[stacks-perfect-section-phantom]{Catégories dérivées de champs}
\item \hyperref[stacks-introduction-section-phantom]{Introduction aux champs algébriques}
\item \hyperref[stacks-more-morphisms-section-phantom]{Compléments sur les morphismes de champs}
\item \hyperref[stacks-geometry-section-phantom]{La géométrie des champs}
\end{enumerate}
Thèmes de théorie des espaces de modules
\begin{enumerate}
\setcounter{enumi}{107}
\item \hyperref[moduli-section-phantom]{Champs de modules}
\item \hyperref[moduli-curves-section-phantom]{Espaces de modules de courbes}
\end{enumerate}
Divers
\begin{enumerate}
\setcounter{enumi}{109}
\item \hyperref[examples-section-phantom]{Exemples}
\item \hyperref[exercises-section-phantom]{Exercices}
\item \hyperref[guide-section-phantom]{Guide bibliographique}
\item \hyperref[desirables-section-phantom]{Desiderata}
\item \hyperref[coding-section-phantom]{Style de programmation}
\item \hyperref[obsolete-section-phantom]{Obsolète}
\item \hyperref[fdl-section-phantom]{Licence de documentation libre GNU}
\item \hyperref[index-section-phantom]{Index généré automatiquement}
\end{enumerate}
\end{multicols}


\bibliography{my}
\bibliographystyle{amsalpha}

\end{document}
